\documentclass[12pt]{article}
\usepackage[utf8]{inputenc}
\usepackage[T1]{fontenc}
\usepackage[compress]{cite}
\bibliographystyle{naturemagnourl}
\usepackage{IEEEtrantools,stackengine}
\stackMath
\usepackage{authblk}
\usepackage{fullpage}
\usepackage{amssymb,amsmath}
\usepackage{adjustbox}
\usepackage{graphicx}
\usepackage{ragged2e}
\usepackage{pdflscape}   
\usepackage{geometry}    
\usepackage{booktabs}    
\usepackage{array}     
\usepackage{siunitx}
  \sisetup{
    group-separator = {,},
    group-minimum-digits = 4,
    detect-all,
    detect-weight=true,
    detect-family=true,
    mode=text
  }
\usepackage{xspace}
\newcommand\ie{i.\,e.\xspace}
\newcommand\eg{e.\,g.\xspace}
\newcommand\Eg{E.\,g.\xspace}
\newcommand\NB{N.\,B.\xspace}
\newcommand\BSc{B.\,Sc.\xspace}
\newcommand\MSc{M.\,Sc.\xspace}
\newcommand\PhD{Ph.\,D.\xspace}
\newcommand\etc{etc.\xspace}
\newcommand\resp{resp.\xspace}
\newcommand\cf{cf.\xspace}
\newcommand\Cf{Cf.\xspace}
\newcommand\cp{cp.\xspace}
\newcommand\etal{et\,al.\xspace}
\newcommand\page[1]{p.\,#1}
\newcommand\pages[1]{pp.\,#1}
\newcommand\pa{p.\,a.\xspace}
\newcommand\ham{a.\,m.\xspace}
\newcommand\hpm{p.\,m.\xspace}
\newcommand\UK{U.\,K.\xspace}
\newcommand\US{U.\,S.\xspace}
\newcommand\wlogenerality{w.\,l.\,o.\,g.\xspace}
\DeclareMathOperator{\logit}{logit}
\newcommand\abs[1]{\left| #1 \right|}
\newcommand{\norm}[1]{\left\lVert#1\right\rVert}
\usepackage{booktabs}
\usepackage{longtable}
\usepackage{multirow}
\usepackage[dvipsnames, table]{xcolor}
\usepackage[left]{lineno}
\renewcommand\linenumberfont{\normalfont\tiny}
\usepackage{tikz}
\usetikzlibrary{arrows.meta, positioning}
\usepackage{minted}
\usepackage{csquotes}

\usepackage{pifont}
\newcommand{\checked}{\ding{110}}
\newcommand{\unchecked}{\ding{111}}

\newcommand{\TODO}[1]{{\color{red}#1}}
\newcommand{\MOVE}[1]{{\color{orange}#1}}
\newcommand{\INSERT}[1]{{\color{blue}#1}}
\definecolor{darkgreen}{rgb}{0.0, 0.5, 0.0}
\newcommand{\FIX}[1]{{\color{darkgreen}#1}}
\definecolor{lightyellow}{HTML}{FFE699}
\definecolor{red_revision}{HTML}{FF0000}

\usepackage{setspace}
\doublespacing
\usepackage[unicode=true]{hyperref}
\hypersetup{breaklinks=true,
%            bookmarks=true,
%            colorlinks=true,
            pdfborder={0 0 0},
            allcolors=black}
\urlstyle{same} % don't use a different (monospace) font for urls
\usepackage[%
  % capitalise,
  sort&compress
]{cleveref}
  \Crefname{appendix}{Supplement}{Supplements}
  \Crefname{figure}{Fig.}{Fig.}
  \crefname{table}{Table}{Tables}   % mid-sentence
  \Crefname{table}{Table}{Tables}   % sentence start
\usepackage{times}
\usepackage{enumerate}
\usepackage{enumitem}
\usepackage{float}
\setcounter{secnumdepth}{5}
% Figure and Tables related;
\usepackage{subfig}
\usepackage{graphicx}
% Plots path
%\graphicspath{{./Plots/}}
\newcommand{\figletter}[1]{{{\fontfamily{\sfdefault}\selectfont \textbf{#1}}}}
\usepackage{array}
\usepackage{tabularx}
\usepackage{dcolumn}
\usepackage{rotating}
% from: http://tex.stackexchange.com/questions/2441/how-to-add-a-forced-line-break-inside-a-table-cell
\newcommand{\mcell}[2][c]{%
  \begin{tabular}[c]{@{}#1@{}}#2\end{tabular}}
\newcommand{\mcellt}[2][c]{%
  \begin{tabular}[t]{@{}#1@{}}#2\end{tabular}}
\newcommand{\lcellt}[2][l]{%
  \begin{tabular}[t]{@{}#1@{}}#2\end{tabular}}
\newcommand{\stackedcell}[2][c]{%
  \begin{tabular}[#1]{@{}c@{}}#2\end{tabular}}
\usepackage{titlesec}

% Redefining subsubsection format
\titleformat{\subsubsection}
  {\normalfont\itshape}{\thesubsubsection}{1em}{}
\titleformat{\subsection}
  {\normalfont\bfseries}{\thesubsection}{1em}{}
\makeatletter
\renewcommand{\fps@figure}{H}         % default {tbp}
\renewcommand{\fps@table}{H}         % default {tbp}
\makeatother 
\renewcommand{\arraystretch}{1.2}
\allowdisplaybreaks
%\hyphenation{CascadeLSTM LSTM}

\usepackage{eurosym}
\usepackage{comment}

\newcommand{\REV}[1]{\textcolor{blue}{#1}}

% Hypothesis
\usepackage{ntheorem}
\theoremseparator{:}
\newtheorem{hyp}{Hypothesis}
\newtheorem*{hyp_non}{Hypothesis}
\usepackage{makecell}

\begin{document}

%%%%%%%%%%%%%%%%%%%%%%%%%%%%%%%%%%%%%%%%%%%%%%%%%%%%%%%%%%%%%%%
% Title page

\title{\centering\LARGE\singlespacing A large-scale systematic map of anthropogenic drivers of biodiversity loss}

\renewcommand\Affilfont{\fontsize{9}{10.8}\selectfont}

% Author names
\author[1]{Hunain Mohuiddin}
\author[1,2]{Kerstin Forster}
\author[,1,2]{Stefan Feuerriegel\thanks{Corresponding author: feuerriegel@lmu.de}}

%\equalcont{These authors contributed equally to this work.}
% \thanks{Correspondence: feuerriegel@lmu.de}

\affil[1]{LMU Munich, Munich, Germany}
\affil[2]{Munich Center for Machine Learning, Munich, Germany}
\date{}
\maketitle
\newpage
\linenumbers

\begin{abstract}\normalfont
\noindent
Biological diversity underpins ecosystem functioning and planetary resilience. Biodiversity research has expanded rapidly, yet keeping track of which anthropogenic drivers, ecosystems, and taxa are most studied remains a major challenge. Here we present a large language model (LLM)--based framework to systematically map the evidence base on anthropogenic drivers of biodiversity loss at scale. Our framework classifies scientific publications by the anthropogenic drivers, specific threats, geographic scope, ecosystems, study designs, and taxa involved, drawing on established biodiversity classification schemes; taxonomic identity is resolved against a comprehensive reference covering over two million described species. Our LLM-based mapping yields a systematic map of 253{,}081 scientific articles on biodiversity loss, revealing disparities in research attention across drivers, ecosystems, and geographies. This map offers a foundation for identifying research gaps, tracking progress, and guiding priorities in biodiversity science and conservation.

\end{abstract}
\flushbottom
\maketitle
\thispagestyle{empty}
\sloppy
\raggedbottom
\newpage

\section*{Main}
\label{sec:main}
% P1 - my hook
Achieving the 2050 Vision of living in harmony with nature, set out in the Kunming--Montreal Global Biodiversity Framework (KM-GBF), requires 
transformative changes across economic, social, political, and technological systems \cite{spm_ga_ipbes_2019}. Central to this transformation is understanding what drives biodiversity loss.

% P2 - Motivation: heterogeneity + interdependency
% optional: in-text ciration- % Refer Figure SPM 1 @ Page 9,  Key Message D.8  of spm_nxs_ipbes_2024; Figure ES.1 on page 27 of unep_peace_with_nature_2021
% streck_synergies_2023, unep_peace_with_nature_2021 - to reread
Biodiversity research is highly heterogeneous: it spans an enormous range of taxa, ecosystems, and geographic regions, and employs diverse methodological approaches across exposures as varied as land-use change, pollution, invasive species, climate change, and direct exploitation. Crucially, these drivers do not act in isolation: land-use pressures interact with climate change, overexploitation compounds habitat degradation, and conservation outcomes depend on which combination of drivers is in play \cite{spm_ga_ipbes_2019, spm_nxs_ipbes_2024}. Any synthesis of the global evidence base must grapple with this breadth and interdependency simultaneously.

% P3 -  core argument: gaps and biases
% reread - diaz_ipbes_2015
Yet the existing literature remains structured by pronounced imbalances. First, a small subset of taxa receives the vast majority of research attention while most species remain understudied \cite{troudet2017taxonomic, gorneau2023measuring, bonnet2002taxonomic, marco_changing_trends_biases}. Second, driver coverage is similarly skewed: climate change accounts for nearly 45\% of driver-related publications, while pollution appears in fewer than 5\% \cite{mazor_et_all_mismatch_of_policy}. Third, existing systematic maps and meta-analyses address a single driver, taxon, or region, reflecting these same biases \cite{feeley2017most, moullec_identifying_2021, jaureguiberry_direct_2022, duenas_threat_2021}. A comprehensive synthesis that jointly characterises the evidence base across all major drivers, specific threats, ecosystems, geographic context, and taxa is therefore absent.

% P4 - Our contribution
Here, we present such a map, a structured overview of over $253,081$ scientific articles on the direct anthropogenic drivers of biodiversity loss, retrieved from the Web of Science Core Collection (WOS) and classified across six dimensions: anthropogenic drivers, specific threats, geographic scope, ecosystems, study designs, and taxa.

% P5 - LLM approach - how is it made possible?
Classifying evidence at this scale is infeasible manually.  We therefore automate the screening and coding pipeline using Large Language Models (LLMs) following Collaboration for Environmental Evidence (CEE) systematic map guidelines \cite{CEE_Guidelines_5_1_2022, nunez-mir_automated_2016}. The result is a structured, comprehensive evidence base that captures the full breadth of biodiversity research and can be continuously updated as new literature accumulates.

% P8 - Close by presenting a brief preview of the upcoming results and discussion 
In this article, ....  \TODO{(preview results summary: TBA at last)} 

%%%%%%%%%%%%%%%%%%%%%%%%%%%
\section*{Results}
\label{sec:results}
% < 1200 Words; 
% Core text (Interpretive instead descriptive): Final narrative synthesis for literature growth, size, and analysis/gaps for using 6 sets of labels

\subsection*{Scientific evidence on biodiversity loss is extensive and rapidly expanding}
Our systematic map comprises 253{,}081 scientific articles reporting biodiversity loss linked to direct anthropogenic drivers. Each article is coded across six dimensions: the five IPBES direct drivers, specific IUCN threats, geographic scope, ecosystem context, study design, and taxonomy. Together, these classifications form a structured, multidimensional evidence base in which research attention can be analyzed jointly across drivers, places, ecosystems, methods, and taxa.

The map reveals that the annual number of scientific articles on biodiversity loss increased from 2{,}530 in 2000 to 20{,}520 in 2025, a compound annual growth rate (CAGR) of 8.73\% (see Figure~\ref{figure:time_development}a). Growth was especially pronounced for climate change and severe weather, which increased at 14.75\% annually compared with the overall rate; the category's share of broad threat attributions rose from 5.3\% to 19.3\% (see Figure~\ref{figure:time_development}b,c). This increasing research emphasis is consistent with previous evidence that scientific attention is unevenly distributed across biodiversity-loss drivers and particularly concentrated on climate change \cite{mazor_et_all_mismatch_of_policy}. However, these proportions describe the distribution of scientific attention, not the relative contribution of each driver to biodiversity loss. The systematic map provides the common evidence base for the analyses that follow, which examine how research attention is distributed across drivers, geographies, ecosystems, and socioeconomic contexts.

% figure:screening_process
\begin{figure}[H]
  \centering
  \includegraphics[width=\textwidth]{attachments/combined-time-development.pdf}
  \caption{\textbf{Temporal development and changing composition of the scientific evidence on biodiversity loss:} (\textbf{a}) Annual number of scientific articles for complete publication years from 2000 to 2025. Selected international biodiversity-policy and assessment milestones are shown for context. (\textbf{b}) Annual distribution of broad IUCN threat categories, expressed as the percentage of all threat attributions assigned in each year. (\textbf{c}) Five-year CAGR for the complete biodiversity-loss evidence base and individual broad threat categories.}
  \label{figure:time_development}
\end{figure}



\subsection*{Biodiversity-loss evidence differs systematically across national income groups}

Global assessments identify land- and sea-use change and direct exploitation as the two direct drivers with the largest impacts on nature to date \cite{spm_nxs_ipbes_2024}. The scientific literature, however, documents these drivers unevenly across national income groups. Across 94{,}443 scientific articles reporting observational evidence of biodiversity loss in a specific country, the documented-threat composition of lower-income-country evidence was dominated by agriculture and aquaculture and by biological resource use, whereas higher-income-country evidence gave greater relative emphasis to biological invasions and to climate change and severe weather. The largest contrasts were for agriculture and aquaculture (16.0 percentage points lower in high-income evidence; 95\% bootstrap confidence interval 15.0--17.0) and biological resource use (12.2 points lower; 11.2--13.3), alongside biological invasions (9.1 points higher; 8.3--9.9) and climate change (6.0 points higher; 5.2--6.7) (see Figure~\ref{figure:income_composition}).

These contrasts persisted after directly standardizing the evidence to a common regional and publication-period composition, indicating that they were not explained by differences in where and when the scientific articles were produced: after joint standardization, lower-income-tier evidence remained more oriented toward agriculture (7.0 points) and biological resource use (5.4 points), and higher-income-tier evidence toward biological invasions (5.2 points) and climate change (3.0 points). The pattern held in direction and rank across alternative study-design, income-classification, and weighting specifications (see Supplementary Results). The prominence of agriculture and direct resource use in lower-income-country evidence is consistent with global spatial analyses showing that agriculture and logging are particularly pervasive across tropical regions, while hunting and trapping are geographically widespread threats to mammals and birds \cite{harfoot_using_2021}, and may also reflect the organization of global production, in which consumption in wealthier economies is linked to biodiversity threats in developing producer countries \cite{lenzen_international_2012}. The greater relative representation of biological invasions and climate change in higher-income-country evidence may, in turn, partly reflect persistent geographic inequalities in biodiversity research, with studies disproportionately concentrated in developed countries with greater economic and research capacity \cite{titley_scientific_2017}. These profiles describe documented research attention, not the ecological prevalence or severity of threats.

% figure:income_composition
\begin{figure}[H]
  \centering
  \includegraphics[width=\textwidth]{attachments/income_composition.pdf}
  \caption{\textbf{Documented-threat composition of the biodiversity-loss evidence differs across national income groups:} Threat attributions are drawn from 94{,}443 unique publications reporting negative-direction, observational evidence linked to a country of study, with each publication's attributions allocated fractionally across the broad IUCN threat categories it reports. (\textbf{a}) Within-group composition, expressed as each threat category's share of all attributions in a given World Bank income group; income group is matched to the publication year. (\textbf{b}) High-income-minus-low-income difference in each category's attribution share, in percentage points; points show the observed difference and horizontal lines the 95\% bootstrap confidence interval, ordered by the absolute size of the contrast. Differences shown here are unadjusted; contrasts standardized to a common regional and publication-period composition are reported in the Supplementary Results.}
  \label{figure:income_composition}
\end{figure}


\subsection*{The misalignment between research attention and driver impact is consistent across ecosystem realms}

Global assessments hold that the direct drivers of biodiversity loss differ in importance across the terrestrial, freshwater, and marine realms, with land-use change foremost on land and in freshwater and direct exploitation foremost in marine systems \cite{jaureguiberry_direct_2022, spm_nxs_ipbes_2024}. We examined how the five IPBES direct drivers were distributed across scientific articles within each realm (Figure~\ref{figure:realm_composition}). In every realm, documented research attention departed from this hierarchy in the same direction: pollution drew more attention, and direct exploitation less, than their assessed impact would imply. Pollution held the largest share in both freshwater (57\%) and marine (38\%) evidence, while direct exploitation stayed minor everywhere, despite being the most prevalent documented threat to species worldwide, affecting 72\% of threatened and near-threatened species \cite{maxwell_biodiversity_2016}, and the leading assessed marine driver \cite{spm_nxs_ipbes_2024}. Only on land did the assessed leading driver, land-use change (33\%), also lead in attention.

The freshwater realm shows this most clearly. Its biodiversity loss is attributed to five interacting pressures of comparable standing, including flow alteration and habitat degradation alongside pollution \cite{Reid_freshwater}, yet freshwater attention concentrated on pollution more heavily than in any other realm. Where the real causes are many and diffuse, the evidence narrows onto the one easiest to identify and measure. The gap is also widening: attention to direct exploitation fell in every realm over the study period, most steeply in marine evidence, even as it remained the leading assessed marine driver (Supplementary Figure~\ref{sfig:realm_trends}).

These patterns show that documented research attention does not follow assessed impact, departing from it in the same direction across all three realms. This extends a bias documented across taxa and regions \cite{titley_scientific_2017, troudet2017taxonomic} to the direct drivers of loss themselves, revealing a systematic gap between the drivers most represented in the evidence and those assessed to matter most \cite{diaz_ipbes_2015}.

% figure:realm_composition
\begin{figure}[H]
  \centering
  \includegraphics[width=\textwidth]{attachments/realm_composition.pdf}
  \caption{\textbf{The composition of the biodiversity-loss evidence differs across ecosystem realms:} Each bar shows the share of driver attributions accounted for by the five IPBES direct drivers within a realm, based on the scientific articles assigned to that realm ($n$ shown per realm); each article's attributions are allocated fractionally across the drivers it reports. Pollution holds the largest share in both freshwater (57\%) and marine (38\%) evidence, whereas direct exploitation remains a minor share in every realm despite being the most prevalent assessed threat to species and the leading assessed marine driver. Land-use change leads in attention only in the terrestrial realm (33\%).}
  \label{figure:realm_composition}
\end{figure}


\subsection*{The drivers assessed to matter most are studied through the narrowest taxonomic lens}

Research attention is known to favor vertebrates over invertebrates \cite{titley_scientific_2017, troudet2017taxonomic}. 

Our map reconfirms this skew at scale: across 171,791 biodiversity-loss articles with a resolvable taxonomic group, vertebrates received 34.8\% of documented attention against 4.6\% of accepted described species (7.5 times), whereas invertebrates received 19.7\% against 65.0\% (0.30 times). This skew, however, is not a fixed constant: it depends on which driver of loss is under study. The vertebrate share of attention ranged from 25.9\% under climate change to 50.5\% under direct exploitation, a difference of 24.5 percentage points (95\% bootstrap confidence interval 23.9--25.2). Direct exploitation was the only driver that over-represented vertebrates while under-representing all seven other taxonomic groups; pollution was its mirror image, the only driver to over-represent all four invertebrate, fungal, and microbial groups.

The narrowest lenses sit on the drivers assessed to matter most. Land- and sea-use change and direct exploitation are the two drivers assessed to have the largest impacts on nature to date \cite{spm_nxs_ipbes_2024}. These are also the two whose evidence is most concentrated on vertebrates. Pollution, by contrast, is the most taxonomically inclusive: it draws disproportionate attention relative to its assessed impact across ecosystem realms, as shown above. The gap therefore compounds rather than compensates: where the evidence base is thinnest relative to assessed importance, it is also narrowest taxonomically. Growth has not corrected this. Annual publication volume grew 8.11-fold between 2000 and 2025, yet the vertebrate share of animal-taxon attention moved by only $-0.24$ percentage points, from 62.35\% in 2000--04 to 62.11\% in 2020--25 (95\% CI $-1.51$ to $+0.98$): publication growth has enlarged the evidence base without widening its taxonomic lens.

These patterns extend a taxonomic bias documented in prior work \cite{titley_scientific_2017, troudet2017taxonomic} to the driver-specific structure of the evidence base itself. The map describes documented research attention, not ecological importance, conservation need, or the number of species affected by a driver. It also cannot, on its own, distinguish an attention bias from real differences in the taxonomic composition of what each driver affects, a distinction we return to in the Discussion.

% figure:taxon_driver
% \begin{figure}[H]
%   \centering
%   \includegraphics[width=\textwidth]{attachments/taxon_driver_composition.pdf}
%   \caption{\textbf{Taxonomic attention within each IPBES direct driver:} Colour shows research specialization as the log$_2$ location quotient, how over- or under-represented a taxonomic group is within a driver relative to its share across the whole biodiversity-loss corpus; purple is over-represented, orange under-represented, white proportional. The number in each cell is that group's share of attention within the driver. Direct exploitation is the only driver in which vertebrates alone are over-represented; pollution is the only driver in which the invertebrate, fungal, and microbial rows are all over-represented.}
%   \label{figure:taxon_driver}
% \end{figure}



\newpage
\section*{Discussion}
\label{sec:discussion}
% approx 825 words: synthesis, positioning, limitations, mechanism-vs-impact caveat, implications

The four analyses trace one evidence base along three independent axes of research attention: national income group, ecosystem realm, and taxon. Along each axis, documented attention departs from assessed importance, and does so in a consistent direction: pollution research is disproportionately large relative to its assessed impact, while research on land- and sea-use change and direct exploitation, the two drivers assessed to matter most, is comparatively narrow. Because the map codes these axes jointly rather than in isolation, it also shows that the resulting gaps compound rather than cancel: the drivers whose evidence is thinnest relative to assessed importance are the same drivers whose evidence is most concentrated on a single realm and a single taxon. A synthesis built on any one dimension alone would have missed this stacking; a synthesis built on all six together is what makes it visible.

This closes a gap our Introduction identified directly: existing systematic maps and meta-analyses on biodiversity loss each address a single driver, taxon, or region \cite{feeley2017most, moullec_identifying_2021, jaureguiberry_direct_2022, duenas_threat_2021}, and none jointly characterizes the evidence base across drivers, threats, geography, ecosystems, and taxa at this scale. At 253{,}081 classified articles, our map is, to our knowledge, the largest evidence base of its kind for biodiversity loss, and its value lies less in its size alone than in what joint coding of six dimensions makes visible: compounding gaps that single-dimension reviews cannot detect, because no single dimension shows them. A map of this kind offers a layer of evidence beneath future synthesis and assessment work, one against which the coverage of individual reviews, or the citation practices of assessment bodies, could in principle be checked.

Three limitations qualify how these results should be read. First, the classification pipeline, while validated, is not error-free, and its reliability varies by task: search comprehensiveness reached 85.54\% recall against an a priori test list, screening recall reached 90.45\% on the held-out test split, and coding performance averaged 83.54\% macro-F1 (range 57.52--98.90\%) and Cohen's $\kappa$ of 0.72 (range 0.25--0.93) across the six label sets (see Supplementary Methods for the full validation results). Findings resting on less-validated labels, such as ecosystem sub-typology or species-level taxonomy, carry proportionally more uncertainty than those resting on driver, geographic, or genus-level labels, and should be weighted accordingly (Supplementary Table~\ref{tab:coding_validation_status}). Second, the analyses presented here narrow progressively from the full evidence base: screening identified 605{,}199 eligible records reporting biodiversity change of any direction, of which 253{,}081 reported loss and form the analyzed dataset, and individual findings draw on further subsets, such as the 94{,}443 articles with a documented country of study underlying the income analysis. Third, all compositional and temporal analyses are bounded to complete publication years, 2000 through 2025; typical indexing lag means even the final year is likely a slight undercount, and drivers or taxa with recently accelerating attention are the most susceptible to this.

A further caveat applies specifically to the attention-versus-importance comparisons. Departures from assessed impact are consistent with under-attention, but the map cannot fully separate that explanation from genuine differences in what each driver, realm, or income context actually affects: exploitation-related research may concentrate on vertebrates partly because fisheries and wildlife harvest genuinely target vertebrates more than other groups, and pollution research may draw on invertebrate and microbial model organisms for established ecotoxicological and regulatory reasons rather than because attention is unusually inclusive. This ambiguity is not unique to this map; it is inherent to any evidence-attention analysis that lacks an independent benchmark for what a driver actually affects, and it also applies to prior work documenting taxonomic bias \cite{titley_scientific_2017, troudet2017taxonomic}. Only the taxonomic anchor offers such a benchmark here, against described species richness, and even that benchmark is suggestive rather than conclusive, since described diversity is not the same as what a specific driver targets. Testing the alternative directly would require comparable external benchmarks for the other axes, such as catch and harvest composition for direct exploitation or standard test-organism conventions for ecotoxicology.

These limitations narrow how strongly the findings should be read, not what they are useful for. The map points to where new evidence would close the largest documented gaps relative to assessed importance: non-vertebrate impacts of direct exploitation and land- and sea-use change, and driver evidence beyond pollution in freshwater and marine systems, whose actual pressures are documented as more numerous and interacting than the literature reflects \cite{Reid_freshwater}. It also distinguishes evidence that is already rich but has not been synthesized from evidence that remains genuinely thin, a distinction future systematic reviews could use to prioritize where synthesis, rather than further primary research, is the more efficient next step. More broadly, a jointly coded map of this kind offers a check available to assessment bodies such as IPBES: whether the evidence cited in an assessment tracks the underlying distribution of research attention, or draws disproportionately on its most convenient or most-cited subset. Resolving the attention-versus-importance ambiguity raised above, through external benchmarks for exploitation and pollution targets analogous to the taxonomic anchor used here, is the most direct next step this map motivates.


\newpage
\section*{Methods}
\subsection*{Systematic mapping protocol}
We used an LLM to create a comprehensive systematic map of direct anthropogenic drivers of biodiversity loss, following CEE guidelines for systematic mapping~\cite{CEE_Guidelines_5_1_2022} with adjustments for our LLM-based pipeline. All steps are documented in a protocol pre-registered prior to full-corpus application. Figure~\TODO{X} summarizes the five-stage workflow, comprising record retrieval, label acquisition, prompt definition for screening, prompt definition for coding, and full-corpus application, each detailed in the following subsections. LLM-specific configurations and rationales for each stage are documented in an accompanying reporting checklist~\cite{feuerriegel_reporting_2026} (Extended Data Table~\ref{tab:guide_llm_checklist}).

\subsection*{Document search}
The first stage retrieved bibliographic records from the WOS. We constructed a search string designed to prioritize recall, combining driver-specific terms adopted from a prior evidence synthesis~\cite{jaureguiberry_direct_2022} with a broad set of biodiversity-relevant subject categories~\cite{moullec_identifying_2021} (see Supplementary Table~\ref{tab:search_string}). Comprehensiveness was assessed against an \emph{a priori} test list of 581 records drawn from five published evidence syntheses~\cite{moullec_identifying_2021, jaureguiberry_direct_2022, beckmann_conventional_2019, benitez-lopez_impact_2017, urban_accelerating_2015} and the IPBES Global Assessment chapters on direct drivers~\cite{balvanera_chapter_2019}, yielding an overall recall of 85.54\% (see Supplementary Table~\ref{tab:search_recall}). Applied to publications from 2000 to April 2026, the search retrieved approximately 2.4~million bibliographic records after deduplication.

\subsection*{Document screening through LLM}
The second stage screened the retrieved records for evidence of a biodiversity change linked to a direct anthropogenic driver. We implemented screening as a zero-shot classification task using GPT-5-nano (version \texttt{2025-08-07}): each record's title and abstract was assessed through a four-question procedure (see Figure~\ref{figure:screening_process}), and a record was classified as eligible only when all four questions returned \texttt{Yes} or \texttt{Unclear}. To assess reliability, we compared LLM screening decisions against expert judgments from previously published, curated evidence syntheses~\cite{shawetal, kecketal, murphyetal, jaureguiberry_direct_2022, iucn_rle_2022}. A human baseline, established by having a person screen from abstracts alone ($n = 100$), achieved a recall of 83\%. The LLM, evaluated against the full-text expert labels ($n = 2{,}418$, partitioned into training, development, and held-out test splits), achieved a recall of 90.45\% on the test split ($n = 482$). A further check compared LLM and human screening decisions on 400 records drawn directly from search retrieval rather than this curated benchmark, finding a recall of 90.1\% (precision 64.5\%, F1 75.2\%), consistent with the LLM's performance on the curated data. Applied to the full corpus, screening yielded 605{,}199 eligible records, of which 253{,}081 reported a negative direction of change and constitute the biodiversity loss dataset analyzed in this study (Supplementary Figure~\ref{fig:roses}).

\subsection*{Document coding through LLM}
The third stage coded eligible records across six label sets (\texttt{L1}--\texttt{L6}) that together reconstruct an evidence chain from direct anthropogenic driver to biodiversity impact: driver category (\texttt{L1}), specific threats (\texttt{L2}), geographic scope (\texttt{L3}), ecosystem typology (\texttt{L4}), study attributes (\texttt{L5}), and taxonomy (\texttt{L6})~\cite{diaz_ipbes_2015, iucn_threats_classification_2022, data_ipbes_geo, keith2022function_get, GBIFBackboneTaxonomy_general}. Label definitions and decision rules were adopted verbatim from these established schemes; \texttt{L5} follows a custom schema developed for this map. Hierarchical label sets (\texttt{L2}, \texttt{L4}) required multiple LLM calls per record to resolve their threat or ecosystem hierarchy; all others used a single call. As with screening, we assessed reliability by comparing LLM coding decisions against expert judgments: a human baseline established how well a person could code from abstracts alone, and a second check evaluated the LLM against full-text expert labels across all six label sets. On the held-out test split, the pipeline achieved a macro-average weighted F1 of 83.54\% (range: 57.52\%--98.90\%) and Cohen's~$\kappa$ of 0.72 (range: 0.25--0.93).

\subsection*{Validation}
Both screening and coding used GPT-5-nano (version \texttt{2025-08-07}) throughout, selected after comparison against larger alternatives where marginal performance gains did not justify cost at the scale of this map. Reasoning effort was the only varied parameter, selected per label set on the development split. Both reference datasets were partitioned into training, development, and test splits; prompts were developed on the training split, reasoning effort selected on the development split, and the test split held out for final evaluation (see Supplementary Methods).

\newpage
\section*{Data availability}
Data for all $253,081$ articles on biodiversity loss coded in this systematic map are available via Open Science Framework (OSF) at \url{https://osf.io/xg8yq/overview}.  WoS-derived metadata is excluded in compliance with Clarivate's end-user terms; full details on record identifiers, licensing restrictions, and reproduction instructions are documented within the OSF repository.
\TODO{Review once before submission: Add final .csv data in repo along with instructions}
\vspace{0.4cm}

% We provide two categories of data. For the validation datasets (reference and annotated screening and coding datasets; see Supplementary Methods), we provide the DOI, WoS Accession Number (UT), raw extraction queries, and all assigned labels for each record, enabling independent reproduction of the reported validation metrics. For the literature search corpus, retrieved under an institutional WoS license, we provide a study-specific surrogate key and the corresponding \texttt{L1}--\texttt{L6} labels for every record; WoS-derived metadata is excluded, as it cannot be redistributed under Clarivate's end-user terms, though the corpus itself can be reconstructed from the search string in Supplementary Table~\ref{tab:search_string}. Record-level identifiers for independent audit are available from the corresponding author upon request. All data are available at the code repository.
\vspace{0.4cm}

\section*{Code availability}
All code used for the LLM-based screening and coding pipeline and validation is available at \url{https://github.com/hunaindot/biodiversity-taxonomy-thesis/tree/main}. The repository includes prompts, configuration files, and a README detailing setup and reproduction steps.

\vspace{0.4cm}

\newpage
\bibliography{literature}
\newpage

\section*{Acknowledgments}
We thank the authors of the published evidence syntheses used to construct the search test list and the validation datasets for screening and coding of this systematic map
\vspace{0.4cm}

\section*{Author contributions}
H.M. designed the study, developed the LLM-based screening and coding pipeline, performed the analysis, and wrote the manuscript. K.F. reviewed the manuscript and provided feedback. S.F. conceived the study, provided supervisory guidance, and reviewed the manuscript. All authors read and approved the final manuscript.
\vspace{0.4cm}

\section*{Funding} 
No external funding was received.
\vspace{0.4cm}

\section*{Competing interests}
The authors declare no competing interests.

\newpage


%%%%%%%%%%%%%%%%%%%%%%%%%%%%%%%%%%%%%%
\appendix
\section*{Supplementary information}
\setcounter{table}{0}
\setcounter{figure}{0}
\renewcommand{\thetable}{S\arabic{table}}
\renewcommand{\thefigure}{S\arabic{figure}}
\renewcommand{\tablename}{Supplementary Table}
\renewcommand{\figurename}{Supplementary Figure}
This section provides supporting material for the results and methods summarized in the main text: Supplementary Tables, Supplementary Figures, and Supplementary Methods detailing the pipeline underlying the main-text results. All materials required to reproduce these methods are available via the registered protocol, code repository, and data repository
\TODO{Review once again before final submisison}




\newpage
\subsection*{Supplementary tables}

\newpage
\subsection*{Supplementary figures}

% figure:screening_process
\begin{figure}[H]
  \centering
  \includegraphics[width=\textwidth]{attachments/screening_process.pdf}
  \caption{\textbf{Eligibility screening procedure:} Each record is screened using four questions (\texttt{Q1}--\texttt{Q4}) applied to its title and abstract: \texttt{Q1} biodiversity change (does the record report at least one form of biodiversity change, operationalized across genetic, species, community, and ecosystem change); \texttt{Q2} direction of change (does the record identify a direction of biodiversity change, recorded as negative, positive, mixed, or unclear, with negative change, i.e.\ loss, as the primary analytical focus of the findings); \texttt{Q3} direct anthropogenic driver (does the record mention or plausibly imply at least one driver, assessed using the five IPBES driver categories and the IUCN Threats Classification Scheme); and \texttt{Q4} linkage (does the record report or plausibly imply a link between the identified driver and the biodiversity change). A record is classified as eligible only when all four questions return \texttt{Unclear} or \texttt{Yes}; a single \texttt{No} at any question yields an ineligible classification. Each question produces a structured label alongside its result. Full operationalization is provided in the screening prompt in the available code.}
  \label{figure:screening_process}
\end{figure}

\newpage
\begin{figure}[H]
\centering
\includegraphics[width=\textwidth]{attachments/roses.pdf}
\caption{\textbf{ROSES flow diagram:} Of 605{,}199 records meeting the pre-registered eligibility criteria, 253{,}081 reported a negative direction of change and constitute the evidence base for biodiversity loss analyzed in this study. As pre-registered in the protocol, records with a positive, mixed, or unclear direction of change (n = 352{,}118) were retained in the eligible set but excluded from synthesis, since biodiversity loss is the primary analytical focus of this map. A detailed breakdown of records excluded at the eligibility screening stage, by exclusion reason, is provided in Supplementary Figure~\ref{fig:screening_criteria_overlap}.}
\label{fig:roses}
\end{figure}


\newpage
\begin{figure}[H]
\centering
\includegraphics[width=\textwidth]{attachments/screening-criteria-overlap.pdf}
\caption{\textbf{Overlap of screening criteria among ineligible records:} Each column shows one exact combination of screening questions (\texttt{Q1}--\texttt{Q4}, see Supplementary Figure~\ref{figure:screening_process}) answered \texttt{No}, ordered by record count. Most exclusions failed all four criteria simultaneously (n = 1{,}081{,}249, 59.8\% of classified exclusions); a further 24.10\% (n = 435{,}966) failed \texttt{Q1}, \texttt{Q2}, and \texttt{Q4} jointly, and 12.03\% (n = 217{,}696) failed \texttt{Q3} and \texttt{Q4} jointly. The remaining combinations each contributed between 0.01\% and 2\% of classified exclusions.}
\label{fig:screening_criteria_overlap}
\end{figure}


\newpage
% sfig:realm_trends
\begin{figure}[H]
  \centering
  \includegraphics[width=\textwidth]{attachments/realm_trends.pdf}
  \caption{\textbf{Research Attention to pollution and direct exploitation over time within each ecosystem realm, 2000--2025:} Faint lines and points show annual estimates; solid lines show centered five-year means. (\textbf{a}) Pollution as a share of driver attributions within each realm, which remained broadly stable across the period. (\textbf{b}) The share of pollution articles using plastic or marine-debris terminology, which rose from a negligible fraction to roughly a quarter of marine pollution articles and to around one in ten in the terrestrial and freshwater realms. (\textbf{c}) Direct exploitation as a share of driver attributions within each realm, which declined in all three realms and most steeply in marine evidence. Panels (a) and (c) share the fractional-attention definition used in Figure~\ref{figure:realm_composition}.}
  \label{sfig:realm_trends}
\end{figure}


\newpage
\subsection*{Supplementary Results}

\newpage
\subsubsection*{Robustness of the income-composition contrast}
\phantomsection\label{sresults:income_composition}

This section documents the analytical choices and robustness checks underlying the finding that the documented-threat composition of the biodiversity-loss evidence differs across national income groups. The unit of analysis is a scientific article reporting negative-direction, observational evidence for a country of study; articles reporting more than one broad IUCN threat category contribute fractionally, each of the $k$ reported categories receiving weight $1/k$.

\paragraph*{Historical income matching:} Because a country's World Bank income classification can change over time, each article was matched to the income group in force in its publication year rather than to the current classification. For complete publication years 2000--2025, this yielded 94{,}443 articles spanning 213 countries, and reclassified 11.0\% of article--country assignments relative to current income groups. The composition contrasts below are therefore computed on income groups as they stood at the time of study.

\paragraph*{Standardization for region and period:} Income groups differ in which world regions and publication periods they draw on, so the raw contrasts could reflect the geography and timing of research rather than income. To test this, we grouped the four income groups into a lower-income tier (low and lower-middle) and a higher-income tier (upper-middle and high), then directly standardized both tiers to a common distribution over World Bank region and five-year publication period, using the 28 region--period strata present in both tiers. The four principal contrasts persisted (Table~\ref{tab:income_standardized}): after joint standardization, lower-income evidence remained more oriented toward agriculture and biological resource use, and higher-income evidence toward biological invasions and climate change, with all four bootstrap intervals excluding zero. Pollution behaved differently: its small raw contrast fell to near zero under region adjustment, so we do not treat it as a stable income difference.

% tab:income_standardized
\begin{table}[H]
\renewcommand{\arraystretch}{1.3}
\setlength{\tabcolsep}{6pt}
\centering
\small
\begin{tabularx}{\linewidth}{|>{\raggedright\arraybackslash}X
                             |>{\centering\arraybackslash}p{0.11\linewidth}
                             |>{\centering\arraybackslash}p{0.13\linewidth}
                             |>{\centering\arraybackslash}p{0.11\linewidth}
                             |>{\centering\arraybackslash}p{0.15\linewidth}|}
\hline
\textbf{Threat category} & \textbf{Raw} & \textbf{Region-std.} & \textbf{Joint-std.} & \textbf{95\% CI (joint)} \\
\hline
Invasive \& problematic species & 6.2 & 5.3 & 5.2 & 4.4 to 6.0 \\
\hline
Climate change \& severe weather & 3.2 & 3.9 & 3.0 & 1.1 to 3.9 \\
\hline
Pollution & 1.1 & $-0.1$ & 1.5 & 0.2 to 2.8 \\
\hline
Natural system modifications & 3.3 & 0.8 & 1.2 & 0.5 to 1.8 \\
\hline
Biological resource use & $-7.0$ & $-5.6$ & $-5.4$ & $-6.3$ to $-4.4$ \\
\hline
Agriculture \& aquaculture & $-6.5$ & $-6.0$ & $-7.0$ & $-7.9$ to $-5.9$ \\
\hline
\end{tabularx}
\caption{\textbf{Higher-income-minus-lower-income threat-share contrasts under standardization:} Differences in each broad IUCN threat category's share of attributions between the higher- and lower-income tiers, in percentage points, shown raw, after region-only direct standardization, and after joint region-and-period standardization. Positive values indicate greater relative attention in higher-income evidence. Confidence intervals are bootstrap intervals for the joint-standardized contrast. Categories with contrasts near zero across all specifications are omitted.}
\label{tab:income_standardized}
\end{table}

\paragraph*{Sensitivity analyses:} The pattern was stable across alternative specifications. Using all negative-direction study designs rather than observational articles only, using current rather than historical income classification, shifting the fiscal-year boundary, including the partial 2026 year, and applying country-fractional weighting each preserved the direction of every contrast and the rank ordering of threats (Spearman rank correlation 0.99 to 1.00), with no single contrast shifting by more than 4.2 percentage points. These results describe documented research attention, not the ecological prevalence or severity of threats.




\newpage
\subsection*{Supplementary methods}
This section provides the methodological details underlying the methods section, covering the literature search, validation datasets, and screening and coding results. 

\newpage
\subsubsection*{Supplementary methods: Search}
\phantomsection\label{smethods:search}
The literature search was conducted on the WOS for publications from 2000 to April 2026, with no language restrictions. The search string was designed to prioritize recall over precision; its five components are combined with the \texttt{AND} operator (see Supplementary Table~\ref{tab:search_string}). Component~(A) restricts retrieval to journal articles; component~(B) limits results to biodiversity-relevant WoS subject categories; component~(C) requires language indicating measured change, assessment, or monitoring; component~(D) supplies five driver-specific sub-strings, one per IPBES direct driver category, adopted from a prior rapid evidence assessment~\cite{jaureguiberry_direct_2022}, combined with \texttt{OR}; and component~(E) applies the publication-year filter. Gray literature was excluded because its processing at this corpus scale falls outside the scope of the LLM pipeline, a recognized limitation of the search coverage.

We assessed search comprehensiveness against an \emph{a priori} test list of 581 records drawn from five published sources (three meta-analyses, an evidence assessment, and a systematic map protocol) on direct anthropogenic drivers of biodiversity change~\cite{moullec_identifying_2021, jaureguiberry_direct_2022, beckmann_conventional_2019, benitez-lopez_impact_2017, urban_accelerating_2015} and the IPBES Global Assessment chapters on direct drivers (sections 2.1.12--2.1.17)~\cite{balvanera_chapter_2019}. The list was assembled by extracting reference lists from each source and excluding gray literature and records not indexed in WoS, yielding 581 records from an initial 791 candidates; scientific publications cited by the original authors as evidence were treated as eligible for this map. A comprehensiveness check conducted on 02~May~2026 found that the final search string retrieved 497, an overall recall of 85.54\%, with source-specific recall ranging from 71.43\% to 97.32\% (see Supplementary Table~\ref{tab:search_recall}).

% tab:search_string
\newgeometry{top=0.8cm, bottom=1.5cm, left=1cm, right=1cm}
\thispagestyle{plain}
\begin{landscape}
\begin{table}[p]
\renewcommand{\arraystretch}{1.15}
\setlength{\tabcolsep}{3pt}
\centering
\tiny
\begin{tabularx}{\linewidth}{|>{\raggedright\arraybackslash}p{0.13\linewidth}
                             |>{\raggedright\arraybackslash}p{0.06\linewidth}
                             |>{\raggedright\arraybackslash}p{0.22\linewidth}
                             |>{\raggedright\arraybackslash}p{0.06\linewidth}
                             |>{\raggedright\arraybackslash}X|}
\hline
\textbf{Component} &
\textbf{Operator} &
\textbf{What this block does} &
\textbf{WoS field} &
\textbf{Search string} \\
\hline

(A) Publications &
AND &
Limit results to articles only.  &
DT &
\texttt{DT=(Article)} \\
\hline

(B) Scientific disciplines &
AND &
Limits results to relevant Web of Science Categories to keep retrieval manageable. &
WC &
\texttt{WC=("Biodiversity Conservation" OR "Ecology" OR "Environmental Sciences" OR "Environmental Studies" OR "Marine \& Freshwater Biology" OR "Fisheries" OR "Oceanography" OR "Forestry" OR "Water Resources" OR "Agronomy" OR "Geosciences, Multidisciplinary" OR "Meteorology \& Atmospheric Sciences" OR "Plant Sciences" OR "Zoology" OR "Remote Sensing" OR "Multidisciplinary Sciences")} \\
\hline

(C) Change / assessment / monitoring terms &
AND &
Ensures retrieved records include language indicating measured change or assessment/monitoring (e.g., loss/change, impacts/pressures, trends/responses, risk, monitoring). &
TS &
\texttt{TS=((loss* NEAR/3 (habitat* OR biodiversit* OR species OR population* OR "ecosystem service*" OR ecosystem* OR forest* OR coral* OR mangrove* OR reef* OR wetland* OR bird* OR mammal* OR plant*)) OR change* OR reference* OR indicator* OR impact* OR effect* OR pressure* OR trend* OR response* OR status OR pattern* OR assessment* OR monitor* OR risk* OR driver*)} \\
\hline

\multicolumn{5}{|l|}{\textbf{(D) Direct drivers of anthropogenic biodiversity change (driver sub-strings combined with OR within TS)}} \\
\hline

(D1) Land/sea use change &
OR &
Terms capturing land/sea use change and habitat modification. &
TS &
\texttt{(("land use chang*" OR "land cover chang*" OR "land* system* chang*" OR "land* chang*" OR "land* degradation" OR "land* conversion" OR "land management") OR ((marine OR ocean OR sea*) AND use AND change*) OR ("habitat loss*" OR "habitat degradation*" OR "habitat chang*" OR "habitat conversion" OR "habitat modification*" OR "habitat fragmentation" OR "habitat management") OR (biofuel* OR biogas*) OR ((residential OR urban* OR cities OR housing OR commercial OR industrial OR touris* OR recreation*) AND (extent OR area OR cover* OR proportion* OR percentage*) AND (develop* OR expan* OR increas*)) OR (deforestation OR afforestation OR "wood plantation*" OR "forest* plantation*" OR "pulp* plantation*") OR ((aquaculture* OR agricultur* OR crop* OR "cultivated area*" OR farmland OR pasture) AND (develop* OR expan* OR increas* OR intensification OR extensification OR abandon* OR chang*)) OR (livestock AND (farm* OR ranch* OR graz* OR densit*)) OR (energy AND (drill* OR min* OR quarry*)) OR (((transport* OR service) AND (network* OR corridor* OR infrastructure*)) OR road* OR rail* OR "ship* lane*" OR "flight path*") OR ("system* modification*" OR (dams OR (water AND (management OR use))))) } \\
\hline

(D2) Direct Exploitation and Resource Extraction &
OR &
Terms capturing extraction/harvest/withdrawal of biological and abiotic resources, including fishing and logging. &
TS &
\texttt{(((((natur* AND resource*) OR material* OR (biomass OR "fossil fuel*" OR ore* OR mineral*)) OR ((biological AND resource*) OR wildlife OR animal* OR plant* OR population*) OR water) AND ((overextraction OR overexploitation) OR (extraction OR exploitation OR withdraw*) OR depletion OR (hunting OR collect* OR gather* OR poaching OR "bush meat*"))) OR ((forest* OR wood OR timber*) AND (harvest* OR logging*)) OR (soil* AND (erosion OR acidification OR degradation)) OR (fisher* OR fishing OR (bycatch OR by-catch))} \\
\hline

(D3) Climate change &
OR &
Terms capturing climate change and associated climatic stressors and impacts. &
TS &
\texttt{("climat* change" OR "chang* climat*" OR ("increase* temperature*" OR "thermal alteration" OR "global warming") OR (chang* AND precipitation*) OR ((extreme OR severe) AND (weather OR temperature OR precipitation OR climat*)) OR heatwave* OR drought* OR storm* OR flood* OR (UV OR ultraviolet) OR "sea level rise" OR "soil salinization" OR ("habitat shift*" AND climate) OR ("habitat alteration" AND climate) OR ((consumptive OR human OR agricultural) AND "water stress"))} \\
\hline

(D4) Pollution &
OR &
Terms capturing pollution and contaminants across environmental media (soil/water/air). &
TS &
\texttt{((pollut* OR contamin* OR emission* OR spill-over* OR disposal* OR deposition* OR load* OR dump* OR discharge*) OR ((noise OR light OR gas* OR particle* OR particulate* OR nitrogen OR sulphur* OR phosphor* OR waste* OR garbage OR sewage OR pesticide* OR fertilizer* OR nutrient* OR "oil spill*" OR metal OR salinization OR salinisation OR acidification OR solid OR plastic OR mercury OR nutrient* OR pollutant*) AND (soil OR water OR ocean OR marine OR atmosphere OR air)) OR ((industrial OR military OR agricultur* OR forest*) AND effluent*))} \\
\hline

(D5) Invasive alien species &
OR &
Terms capturing invasive alien species introductions and biological invasions. &
TS &
\texttt{(((invasive OR invasion) AND (alien OR exotic OR non-native OR pest* OR disease) AND species) OR "biological invasion*" OR (introduc* AND (species OR material)))} \\
\hline

(E) Publication year &
AND &
Restricts results to the selected publication year range. &
PY &
\texttt{PY=2000--2026} \\

\hline

\end{tabularx}
\caption{\textbf{Search string components used in Web of Science}: Components (A)--(E) are combined using \texttt{AND}. Within \texttt{TS}, the five direct driver sub-strings (D1--D5) from \cite{jaureguiberry_direct_2022} are combined using \texttt{OR}. Component~(E) applies the WoS publication-year filter \texttt{PY=2000--2026}; because WoS publication-year filtering does not express a day/month cutoff, records with publication-date metadata after 30~April~2026 were excluded at the point of export.}
\label{tab:search_string}
\end{table}
\end{landscape}
\restoregeometry


% tab:search_recall
\begin{table}[H]
\renewcommand{\arraystretch}{1.3}
\setlength{\tabcolsep}{4pt}
\centering
\scriptsize
\begin{tabularx}{\linewidth}{|>{\raggedright\arraybackslash}X
                             |>{\raggedright\arraybackslash}p{0.19\linewidth}
                             |>{\centering\arraybackslash}p{0.14\linewidth}
                             |>{\centering\arraybackslash}p{0.09\linewidth}
                             |>{\centering\arraybackslash}p{0.09\linewidth}|}
\hline
\textbf{Source} & \textbf{DOI} & \textbf{Test list ($n$)} & \textbf{Retrieved} & \textbf{Recall (\%)} \\
\hline
Beckmann \etal~\cite{beckmann_conventional_2019} & 10.1111/gcb.14606 & 105 & 75 & 71.43 \\
\hline
Ben\'itez-L\'opez \etal~\cite{benitez-lopez_impact_2017} & 10.1126/science.aaj1891 & 60 & 52 & 86.67 \\
\hline
Balvanera \etal~(IPBES 2.1.12--2.1.17)~\cite{balvanera_chapter_2019} & 10.5281/zenodo.3831673 & 119 & 96 & 80.67 \\
\hline
Jaureguiberry \etal~\cite{jaureguiberry_direct_2022} & 10.1126/sciadv.abm9982 & 139 & 128 & 92.09 \\
\hline
Moullec \etal~\cite{moullec_identifying_2021} & 10.1186/s13750-021-00234-y & 46 & 37 & 80.43 \\
\hline
Urban~\cite{urban_accelerating_2015} & 10.1126/science.aaa4984 & 112 & 109 & 97.32 \\
\hline
\textbf{Total} & -- & \textbf{581} & \textbf{497} & \textbf{85.54} \\
\hline
\end{tabularx}
\caption{\textbf{Comprehensiveness of the search against the \emph{a priori} test list}: For each source, the test list count is the number of eligible records indexed in WoS, and recall is the percentage retrieved by the final search string. The test list is the union of all sources ($n = 581$), against which overall recall is assessed.}
\label{tab:search_recall}
\end{table}

\newpage
\subsubsection*{Supplementary methods: Consistency-check datasets}
\phantomsection\label{smethods:cc_datasets}

This map deviates from CEE guidelines by using LLMs rather than human 
reviewers and by screening from abstracts rather than full text; its outputs are therefore validated against expert full-text labels: labels assigned by the original authors of the reference sources through full-text reading. Four datasets support this validation: two drawn from previously published evidence syntheses containing those expert labels, and two manually annotated by the authors of this map from titles and abstracts only. Consistency Check~1 (CC1) compares author abstract-level labels against expert full-text labels to establish a human-performance baseline; Consistency Check~2 (CC2) evaluates LLM pipeline labels against the expert labels. Supplementary Table~\ref{tab:dataset_overview} summarizes the four datasets; Supplementary Table~\ref{tab:source_inventory} lists the secondary sources from which they were compiled. All datasets are shared in the available data for transparency and reproducibility.

% tab:dataset_overview
\begin{table}[H]
\renewcommand{\arraystretch}{1.3}
\setlength{\tabcolsep}{4pt}
\centering
\scriptsize
\begin{tabularx}{\linewidth}{|>{\raggedright\arraybackslash}p{0.22\linewidth}
                             |>{\raggedright\arraybackslash}p{0.12\linewidth}
                             |>{\raggedright\arraybackslash}p{0.30\linewidth}
                             |>{\raggedright\arraybackslash}X|}
\hline
\textbf{Dataset} & \textbf{Size} & \textbf{Details} &
\textbf{Role} \\
\hline
Reference Screening Dataset & $n = 2{,}418$ &
Records included in three meta-analyses, a rapid evidence assessment,
and the Red List of Ecosystems Database. &
CC2: LLM labels evaluated against expert full-text labels to assess
screening recall \\
\hline
Annotated Screening Dataset & $n = 500$ &
Author-assigned screening labels based on abstracts only, samples drawn from the Reference Screening Dataset and search query results. No full-text access allowed during annotation. 
&
CC1A: author abstract-level labels compared against expert full-text labels. CC1B: LLM labels compared against author-annotated labels.
\\
\hline
Reference Coding Dataset & $n = 7{,}844$ &
Records included in three meta-analyses, two systematic reviews, two
systematic maps, a rapid evidence assessment, and one preprint,
carrying expert full-text labels across applicable label sets
\texttt{L1}--\texttt{L6}. &
CC2: LLM labels evaluated against expert full-text labels across all
six label sets (\texttt{L1}--\texttt{L6}) \\
\hline
Annotated Coding Dataset & $n = 600$ (100 per
label set) &
Author-assigned coding labels based on abstracts only, samples drawn
from the Reference Coding Dataset. No full-text access allowed during
annotation. &
CC1: labels assigned from abstracts compared against expert full-text
labels across all six label sets (\texttt{L1}--\texttt{L6}) \\
\hline
\end{tabularx}
\caption{Overview of the four validation datasets supporting the
screening and data-coding consistency checks.}
\label{tab:dataset_overview}
\end{table}


\paragraph*{Reference Screening Dataset ($n = 2{,}418$):}
Compiled from three meta-analyses, a rapid evidence assessment, and the Red List of Ecosystems Database, the Reference Screening Dataset contains 2,418 records. In each source, the original authors read records in full and retained them as eligible; these full-text inclusion decisions are treated as ground-truth labels, with the alignment between source inclusion criteria and the eligibility criteria of this map treated as a working assumption. The dataset is partitioned into train, development, and test splits (60/20/20, stratified by source) and supports CC2.

\paragraph*{Annotated Screening Dataset ($n = 500$):} 
The Annotated Screening Dataset comprises 500 records split into two
sub-samples and independently annotated by the authors from titles and
abstracts only. The \emph{reference dataset sample} ($n = 100$) is drawn from the Reference Screening Dataset and supports CC1A: author abstract-level
labels are compared against expert full-text labels to establish a
human recall baseline. The \emph{search sample} ($n = 400$) is drawn directly from search-query retrieval results and annotated by the authors from titles and abstracts only; it supports CC1B: a blind comparison of LLM screening decisions against author-annotated labels on search results.

\paragraph*{Reference Coding Dataset ($n = 7{,}844$):}
The Reference Coding Dataset was compiled from three meta-analyses, two systematic reviews, two systematic maps, a rapid evidence assessment, and one preprint. Labels were produced by the original authors of each source through full-text coding as part of their respective reviews; these author-assigned labels are treated here as expert ground truth against which LLM coding outputs are evaluated across applicable labels within each of the six label sets (\texttt{L1}--\texttt{L6}).

\paragraph*{Annotated Coding Dataset ($n = 600$):}
The Annotated Coding Dataset comprises of 100 records per label set drawn from the Reference Coding Dataset, with original expert labels stripped before the authors independently re-labeled each record from titles and abstracts only. The dataset supports CC1: author abstract-level labels are compared against expert full-text labels per label set to establish a human-performance baseline.

\paragraph*{Source inventory:}
The datasets and test list draw on 14 previously published reviews, maps, and databases (one pre-print); Supplementary Table~\ref{tab:source_inventory} lists each source with its contributions across the test list and validation datasets. For each source, raw bibliographic and label data were standardized through three steps: records were matched to their WOS bibliographic entries; label values were harmonized to the controlled vocabulary used in this map where source terminologies differed; and records were restructured so that each row corresponds to one record-label-set pair. Full provenance documentation is shared in the available data.

% tab:source_inventory
 \thispagestyle{plain}
\begin{landscape}
\begin{table}[H]
\renewcommand{\arraystretch}{1.2}
\setlength{\tabcolsep}{3pt}
\centering
\scriptsize
\begin{tabularx}{\linewidth}{|>{\raggedright\arraybackslash}p{0.22\linewidth}
                             |>{\raggedright\arraybackslash}p{0.18\linewidth}
                             |>{\centering\arraybackslash}p{0.07\linewidth}
                             |>{\centering\arraybackslash}p{0.08\linewidth}
                             |>{\raggedright\arraybackslash}X|}
\hline
\textbf{Source} & \textbf{Topic / scope} & \textbf{Test list ($n$)} & \textbf{Ref.\ Screening\ ($n$)} & \textbf{Ref.\ Coding (label[$n$])} \\
\hline
Beckmann \etal~\cite{beckmann_conventional_2019} & Land-use intensification and species richness & 105 & - & - \\
\hline
Balvanera \etal~\cite{balvanera_chapter_2019} & IPBES Global Assessment Ch.\ 2.1, sections 2.1.12--2.1.17 & 119 & - & - \\
\hline
Jaureguiberry \etal~\cite{jaureguiberry_direct_2022} & Direct drivers of biodiversity loss & 139 & 128 & region[$128$], realm[$128$], driver[$128$] \\
\hline
Moullec \etal~\cite{moullec_identifying_2021} & Anthropogenic drivers in the North Sea & 46 & - & driver[$1{,}978$], study\_design[$2{,}007$] \\
\hline
Ben\'itez-L\'opez \etal~\cite{benitez-lopez_impact_2017} & Hunting impacts on tropical mammals and birds & 60 & - & region[$48$], country[$48$], realm[$48$], biome[$48$], kingdom[$49$], phylum[$49$], class[$49$], order[$49$], species[$49$] \\
\hline
Urban~\cite{urban_accelerating_2015} & Extinction risk from climate change & 112 & - & region[$103$], sub-region[$42$], country[$14$], kingdom[$100$], phylum[$37$], class[$32$] \\
\hline
Keck \etal~\cite{kecketal} & Global human impact on biodiversity & - & 1{,}711 & study\_design[$2{,}122$], realm[$2{,}122$], driver[$2{,}122$] \\
\hline
Shaw \etal~\cite{shawetal} & Genetic diversity change & - & 418 & - \\
\hline
Murphy and Romanuk~\cite{murphyetal} & Community change and species richness & - & 142 & - \\
\hline
Red List of Ecosystems Database~\cite{IUCNCEM2022RLE} & Ecosystem risk assessments & - & 31 & - \\
\hline
Wright \etal~\cite{Wrightetal} & IUCN threats and status (mustelids) & - & - & region[$608$], sub-region[$534$], country[$513$], threats\_l0[$563$], threats\_l1[$539$], kingdom[$601$], phylum[$601$], class[$601$], order[$601$], genus[$601$], species[$601$] \\
\hline
Lowry \etal~\cite{Lowryetal} & Biological invasions & - & - & region[$1{,}483$], sub-region[$1{,}483$], country[$1{,}477$], realm[$1{,}252$] \\
\hline
Ridley \etal~\cite{ridleyetal} & Threats to species globally & - & - & study\_design[$693$], realm[$693$] \\
\hline
Pulido-Chadid \etal~\cite{Pulidoetal} & Effectiveness of protected areas & - & - & threats\_l0[$73$], threats\_l1[$52$] \\
\hline
\end{tabularx}
\caption{Source inventory of the 14 published sources (one pre-print) used across the test list, Reference Screening Dataset, and Reference Coding Dataset.}
\label{tab:source_inventory}
\end{table}
\end{landscape}
\restoregeometry

\newpage
\subsubsection*{Supplementary methods: Screening}
\phantomsection\label{smethods:screening}
Screening evaluated the corpus of literature retrieved from WoS against the eligibility criteria pre-registered in the protocol, aiming to identify records reporting an aspect of biodiversity change linked to a direct anthropogenic driver (see Supplementary Figure~\ref{figure:screening_process}). 

To assess the reliability of this fully automated procedure, screening decisions were validated through two consistency checks against expert labels drawn from previously published evidence syntheses. Records were labeled independently by an author from titles and abstracts and by the LLM, with reliability assessed primarily via recall, alongside precision and F1 where relevant. This section presents the ROSES flow diagram summarizing record counts from retrieval to final eligible set, followed by the validation results.

\paragraph*{ROSES flow diagram:}
Supplementary Figure~\ref{fig:roses} presents the ROSES flow diagram~\cite{Haddaway2017ROSES} documenting record counts across the screening pipeline. Of 2{,}414{,}191 records entering eligibility screening, 605{,}199 (25.1\%) were retained as eligible. Exclusions were concentrated in a small number of criterion combinations (see Supplementary Figure~\ref{fig:screening_criteria_overlap}). Most exclusions (59.8\%) failed all four criteria simultaneously. A further 24.10\% failed \texttt{Q1} (biodiversity change), \texttt{Q2} (direction of change), and \texttt{Q4} (driver--biodiversity linkage) jointly, while reporting an anthropogenic driver (\texttt{Q3}). Together, these two combinations accounted for 83.90\% of all classified exclusions.

As pre-registered in the Supplementary Text S1 of the protocol, eligibility does not require a negative direction of change: records reporting positive, mixed, or unclear direction were retained as eligible but excluded from synthesis, since biodiversity loss is the primary analytical focus of this map. Of the 352{,}118 eligible records not analyzed in this study, 141{,}267 (23.34\% of eligible records) reported mixed direction, 112{,}684 (18.62\%) reported positive direction, and 98{,}167 (16.22\%) reported unclear direction (Supplementary Table~\ref{tab:direction_breakdown}).

% tab:direction_breakdown
\begin{table}[H]
\renewcommand{\arraystretch}{1.3}
\centering
\small
\begin{tabular}{|l|r|r|}
\hline
\textbf{Direction of change} & \textbf{Records} & \textbf{\% of eligible} \\
\hline
Negative (biodiversity loss) & 253{,}081 & 41.82\% \\
\hline
Mixed & 141{,}267 & 23.34\% \\
\hline
Positive & 112{,}684 & 18.62\% \\
\hline
Unclear & 98{,}167 & 16.22\% \\
\hline
\end{tabular}
\caption{Distribution of direction of change (\texttt{Q2}) across the 605{,}199 eligible records. Negative direction constitutes the biodiversity loss dataset analyzed in this study (Supplementary Figure~\ref{fig:roses}); the remaining categories are retained in the eligible set but excluded from synthesis.}
\label{tab:direction_breakdown}
\end{table}

\paragraph*{Screening validation results:}
The screening prompt operationalizes the four-question procedure as a zero-shot task: GPT-5-nano receives the title and abstract of each record alongside step-by-step guidelines for \texttt{Q1}--\texttt{Q4}, covering biodiversity change, direction of change, direct anthropogenic driver, and driver--biodiversity linkage, returning a structured \texttt{Yes}, \texttt{Unclear}, or \texttt{No} result and a content label for each step. To prioritize recall, the prompt instructs the model to prefer cautious inclusion over strict exclusion, using \texttt{Unclear} rather than \texttt{No} when evidence is weak.

We first established how well a human reviewer can screen records from titles and abstracts alone, without full-text access. Comparing author abstract-level labels against expert full-text labels, this human baseline achieved a recall of 83\% (CC1A). We then evaluated whether the LLM, working under the same abstract-only constraint, could match this baseline: reasoning effort was varied on the development split, the best-performing configuration selected and confirmed on the training split, and evaluated on a held-out test split (CC2). The LLM achieved a recall of 90.45\% on the held-out test split, exceeding the human baseline (see Supplementary Table~\ref{tab:screening_validation}).

\begin{table}[H]
\renewcommand{\arraystretch}{1.3}
\setlength{\tabcolsep}{4pt}
\centering
\scriptsize
\begin{tabular}{|>{\raggedright\arraybackslash}p{0.08\linewidth}
                |>{\raggedright\arraybackslash}p{0.28\linewidth}
                |>{\centering\arraybackslash}p{0.16\linewidth}
                |>{\centering\arraybackslash}p{0.16\linewidth}
                |>{\centering\arraybackslash}p{0.10\linewidth}|}
\hline
\textbf{Check} & \textbf{Comparison} & \textbf{Split [$n$]} &
\textbf{Reasoning effort} & \textbf{Recall (\%)} \\
\hline
CC1A & Author abstract-level labels vs.\ expert full-text labels
    & [$100$] & -- & 83.00 \\
\hline
\multirow{5}{*}{CC2} &
\multirow{5}{*}{\raggedright LLM labels vs.\ expert full-text labels}
  & \multirow{3}{*}{Dev [$484$]} & Low    & 88.42 \\
 & & & Medium & 90.28 \\
 & & & High   & 90.28 \\
\cline{3-5}
 & & Train [$1{,}452$] & Medium & 89.66 \\
\cline{3-5}
 & & Test [$482$]      & Medium & 90.45 \\
\hline
\end{tabular}
\caption{\textbf{Screening validation results (CC1A and CC2).} Recall is the primary metric, prioritized to minimize false exclusions. CC1A establishes a human recall baseline by comparing author abstract-level labels against expert full-text labels. For CC2, reasoning effort was varied on the development split; the medium configuration was selected and confirmed on the training and held-out test splits.}
\label{tab:screening_validation}
\end{table}

% CC2 results, however, are based on the Reference Screening Dataset, which is drawn from prior published reviews where every record was already deemed eligible by the original authors. A reasonable concern is whether the LLM behavior changes when applied to real-world search results rather than this curated benchmark. To check this, we compared LLM and human screening decisions on a separate sample of 400 records drawn directly from search retrieval (CC1B), treating human labels as ground truth. This comparison measures agreement between the two abstract-only labelers, providing a complementary signal to CC2. The LLM achieved a recall of 90.1\% against human labels (precision 64.5\%, F1 75.2\%; see Supplementary Table~\ref{tab:search_sample_confusion}), closely consistent with CC2 test-split recall of 90.45\%. The lower precision reflects the screening design's recall-prioritizing rule of retaining \texttt{Unclear} cases as eligible, leading the LLM to over-include on ambiguous, uncurated records as intended. Together, these results show that the LLM's abstract-level screening behavior is consistent across both the curated reference dataset and uncurated, real-world search results, with no evidence of degraded recall performance when applied beyond the curated benchmark.

CC2 results, however, are based on the Reference Screening Dataset, which is drawn from prior published reviews where every record was already deemed eligible by the original authors. A reasonable concern is whether the LLM behavior changes when applied to real-world search results rather than this curated benchmark. To check this, we compared LLM and human screening decisions on a separate sample of 400 records drawn directly from search retrieval (CC1B), treating human labels as ground truth. This comparison measures agreement between the two abstract-only labelers, providing a complementary signal to CC2. The LLM achieved a recall of 90.1\% against human labels (precision 64.5\%, F1 75.2\%; see Supplementary Table~\ref{tab:search_sample_confusion}), closely consistent with CC2 test-split recall of 90.45\%.

The lower precision reflects the screening design's recall-prioritizing rule of retaining \texttt{Unclear} cases as eligible: a record was treated as eligible unless one of the four screening stages was explicitly coded as No (\texttt{0}), so uncertain stage outputs (\texttt{-1}) remained eligible by default. Of the 60 false positives in CC1B, 17 (28.3\%) contained at least one uncertain stage decision. Reclassifying these conservatively as ineligible increased precision from 64.5\% to 69.3\% but reduced recall from 90.1\% to 80.2\% (F1 74.3\%): the high-recall operating point partly reflects permissive treatment of uncertainty, not the design rule alone.

The remaining 43 false positives (71.7\%) were confident errors. The corresponding human exclusion reasons showed these errors concentrated near the conceptual boundary of the screening definition: the two largest categories were records lacking a clear biodiversity-change endpoint (13/43, 30.2\%) and records where the anthropogenic driver or its linkage to biodiversity change was absent (12/43, 27.9\%). Most residual precision loss therefore traced not to irrelevant retrievals but to abstracts that were environmentally or biodiversity-adjacent while lacking an explicit instance of biodiversity change, an anthropogenic driver, or a clear link between the two. Together, these results show that the LLM's abstract-level screening behavior remains consistent across both the curated reference dataset and uncurated, real-world search results, with recall, the primary metric guiding the screening design, showing no evidence of degradation beyond the curated benchmark; precision proved more sensitive to how borderline and uncertain cases were resolved.


\begin{table}[H]
\renewcommand{\arraystretch}{1.3}
\setlength{\tabcolsep}{6pt}
\centering
\scriptsize
\begin{tabular}{|>{\raggedright\arraybackslash}p{0.22\linewidth}
                |>{\centering\arraybackslash}p{0.14\linewidth}
                |>{\centering\arraybackslash}p{0.14\linewidth}|}
\hline
\multicolumn{1}{|c|}{} & \multicolumn{2}{c|}{\textbf{Human (truth)}} \\
\cline{2-3}
& \textbf{Eligible} & \textbf{Ineligible} \\
\hline
\textbf{LLM Eligible}   & 109 & 60  \\
\hline
\textbf{LLM Ineligible} & 12  & 219 \\
\hline
\end{tabular}
\caption{\textbf{Search-sample agreement matrix (CC1B):} LLM and human abstract-only screening decisions on 400 records drawn directly from search retrieval results, treating human labels as ground truth and "Eligible" as positive class.}
\label{tab:search_sample_confusion}
\end{table}


\paragraph*{Screening output fields:}
Each bibliographic record processed through the screening pipeline carries a structured output comprising three tiers, regardless of eligibility decision. First, provenance fields support full transparency and reproducibility of each screening decision. Second, per-question structured labels capture the reasoning behind each decision: biodiversity change type(s) for \texttt{Q1}, direction of change for \texttt{Q2}, identified driver(s) for \texttt{Q3}, and linkage type for \texttt{Q4}. These labels provide richer evidence than the binary decision alone. Third, a binary eligibility decision (\texttt{pred\_screening}) records whether the record passed the four-question procedure; $605,199$ records were retained as eligible. The complete output schema is presented in Supplementary Table~\ref{tab:screening_output_fields}.

\begin{table}[H]
\renewcommand{\arraystretch}{1.3}
\setlength{\tabcolsep}{4pt}
\centering
\scriptsize
\begin{tabularx}{\linewidth}{|>{\raggedright\arraybackslash}p{0.22\linewidth}|>{\centering\arraybackslash}p{0.12\linewidth}|>{\raggedright\arraybackslash}X|}
\hline
\textbf{Field} & \textbf{Type} & \textbf{Description} \\
\hline
\multicolumn{3}{|l|}{\textit{Metadata}} \\
\hline
\texttt{\_surrogate\_key} & String & Primary key of the record processed in the screening and coding stage. \\
\hline
\texttt{\_source\_run} & String & Identifier for the processing partition in which the record was screened. \\
\hline
\texttt{model} & String & Name of the LLM used to generate the screening output. \\
\hline
\texttt{created\_at} & Datetime & Timestamp indicating when the record was processed. \\
\hline
\texttt{raw\_output} & String & Original LLM response before parsing into structured fields. \\
\hline
\multicolumn{3}{|l|}{\textit{Screening outputs \texttt{Q1}--\texttt{Q4}}} \\
\hline
\texttt{q1\_r} & Integer & Q1 result: biodiversity change reported? $1$ = Yes, $-1$ = Unclear, $0$ = No. \\
\hline
\texttt{q1\_label} & List[String] & Biodiversity change type(s): genetic change, species change, community change, or ecosystem change. \\
\hline
\texttt{q2\_r} & Integer & Q2 result: direction identifiable? $1$ = Yes, $-1$ = Unclear, $0$ = No. \\
\hline
\texttt{q2\_label} & List[String] & Direction label(s): negative, positive, mixed, or unclear. \\
\hline
\texttt{q3\_r} & Integer & Q3 result: direct anthropogenic driver mentioned or plausibly implied? $1$ = Yes, $-1$ = Unclear, $0$ = No. \\
\hline
\texttt{q3\_label} & List[String] & Driver or threat name(s) identified during screening; these support eligibility screening and are not the final coding outputs. \\
\hline
\texttt{q4\_r} & Integer & Q4 result: driver--biodiversity link reported or plausibly implied? $1$ = Yes, $-1$ = Unclear, $0$ = No. \\
\hline
\texttt{q4\_label} & List[String] & Linkage type(s) describing how the driver is linked to the reported biodiversity change. \\
\hline
\multicolumn{3}{|l|}{\textit{Eligibility decision}} \\
\hline
\texttt{pred\_screening} & List[String] & Final eligibility decision derived from \texttt{Q1}--\texttt{Q4}: \texttt{["ELIGIBLE"]} or \texttt{["INELIGIBLE"]}. \\
\hline
\end{tabularx}
\caption{\textbf{Output fields recorded for each record during screening:} Screening fields (\texttt{Q1}--\texttt{Q4}) record the decision result and structured label(s) generated for each eligibility question.}
\label{tab:screening_output_fields}
\end{table}


\newpage

\subsubsection*{Supplementary methods: Coding}
\phantomsection\label{smethods:coding}
All records classified as eligible during screening entered data coding across six label sets (\texttt{L1}--\texttt{L6}), together capturing the evidence chain from anthropogenic driver to biodiversity impact. Prompts, configurations, and the full output schema for each label set are shared in the available code; validation status per label is summarized in Supplementary Table~\ref{tab:coding_validation_status}. Of the 23 labels produced, 14 are validated through full CC1 and CC2 consistency checks and form the primary basis for the results below; the remaining 11 are weak, for reasons detailed later in this section. This section presents CC1 coding validation, followed by CC2 coding validation, and closes with an account of the weak labels.

\newgeometry{top=0.8cm, bottom=1.5cm, left=1cm, right=1cm}
\thispagestyle{plain}
\begin{landscape}
\begin{table}[p]
\renewcommand{\arraystretch}{1.15}
\setlength{\tabcolsep}{4pt}
\centering
\scriptsize
\begin{tabularx}{\linewidth}{|>{\raggedright\arraybackslash}p{0.08\linewidth}
                             |>{\raggedright\arraybackslash}p{0.20\linewidth}
                             |>{\raggedright\arraybackslash}p{0.12\linewidth}
                             |>{\raggedright\arraybackslash}X
                             |>{\centering\arraybackslash}p{0.09\linewidth}
                             |>{\centering\arraybackslash}p{0.09\linewidth}|}
\hline
\textbf{Label set} & \textbf{Label} & \textbf{Type} & \textbf{Description} & \textbf{Is Validated} & \textbf{Is Weak} \\
\hline
\multicolumn{6}{|l|}{\textit{L1: IPBES Direct Drivers}} \\
\hline
L1 & \texttt{driver} & List[String] & IPBES direct driver category or categories assigned to the record. & Yes & - \\
\hline
\multicolumn{6}{|l|}{\textit{L2: IUCN Threats}} \\
\hline
L2 & \texttt{pred\_threat\_l0} & List[String] & Broadest IUCN Threats Classification level assigned to the record. & Yes & - \\
\hline
L2 & \texttt{pred\_threat\_l1} & List[String] & Intermediate IUCN Threats Classification level assigned to the record. & Yes & Yes* \\
\hline
L2 & \texttt{pred\_threat\_l2} & List[String] & Most specific IUCN Threats Classification level assigned to the record. & - & Yes \\
\hline
\multicolumn{6}{|l|}{\textit{L3: Geographic Scope}} \\
\hline
L3 & \texttt{pred\_regions} & List[String] & IPBES region or regions of study focus. & Yes & - \\
\hline
L3 & \texttt{pred\_subregions} & List[String] & IPBES sub-region or sub-regions of study focus. & Yes & - \\
\hline
L3 & \texttt{pred\_countries} & List[String] & Countries of study focus, recorded using ISO 3166-1 alpha-3 codes. & Yes & - \\
\hline
L3 & \texttt{locales} & List[String] & Free-text names of sub-national locations mentioned in the title or abstract, such as parks, reserves, cities, or mountain ranges. & - & Yes \\
\hline
L3 & \texttt{locale\_coordinates} & List[List[Float]] & [latitude, longitude] pairs in decimal degrees (WGS84, EPSG:4326); resolved from explicit coordinates in the title or abstract, or inferred from the locale name if absent. & - & Yes \\
\hline
\multicolumn{6}{|l|}{\textit{L4: Ecosystem Typology}} \\
\hline
L4 & \texttt{pred\_realm} & List[String] & Global Ecosystem Typology realm or realms of study focus. & Yes & - \\
\hline
L4 & \texttt{pred\_biome} & List[String] & Global Ecosystem Typology biome or biomes of study focus. & - & Yes* \\
\hline
L4 & \texttt{pred\_efg} & List[String] & Global Ecosystem Typology Ecosystem Functional Group or groups of study focus. & - & Yes* \\
\hline
\multicolumn{6}{|l|}{\textit{L5: Study Attributes}} \\
\hline
L5 & \texttt{pred\_study\_design} & String & Study design category: observational, experimental, modelling, or unclear. & Yes & - \\
\hline
L5 & \texttt{pred\_methods\_data\_collection} & List[String] & Data collection method or methods reported in the title or abstract, where applicable. & - & Yes \\
\hline
L5 & \texttt{pred\_methods\_analysis} & List[String] & Analysis method or methods reported in the title or abstract, where applicable. & - & Yes \\
\hline
L5 & \texttt{pred\_has\_comparison} & Boolean & True if the record explicitly states or strongly implies a comparator, such as pre/post, control versus exposed sites, or reference sites; otherwise false. & - & Yes \\
\hline
L5 & \texttt{pred\_comparison\_types} & List[String] & Type or types of comparison identified for the record, where applicable. & - & Yes \\
\hline
\multicolumn{6}{|l|}{\textit{L6: Taxonomy}} \\
\hline
L6 & \texttt{kingdom} & List[String] & Kingdom(s) resolved through the GBIF Backbone Taxonomy. & Yes & - \\
\hline
L6 & \texttt{phylum} & List[String] & Phylum/phyla resolved through the GBIF Backbone Taxonomy. & Yes & - \\
\hline
L6 & \texttt{class} & List[String] & Class(es) resolved through the GBIF Backbone Taxonomy. & Yes & - \\
\hline
L6 & \texttt{order} & List[String] & Order(s) resolved through the GBIF Backbone Taxonomy. & Yes & - \\
\hline
L6 & \texttt{genus} & List[String] & Genus/genera resolved through the GBIF Backbone Taxonomy. & Yes & - \\
\hline
L6 & \texttt{species} & List[String] & Species resolved through the GBIF Backbone Taxonomy. & Yes & Yes* \\
\hline
\end{tabularx}
\vspace{4pt}
\caption{\textbf{Coding output fields and validation status:} Of the 23 labels produced during data coding, 14 are designated \emph{validated}: evaluated through the full CC1 and CC2 consistency checks and formed the primary basis for the validation metrics and the results. 11 labels are designated \emph{weak}, for one of three reasons. 7 weak labels (unmarked) were not validated from the outset, as no standardized labeling schema comparable to the classification schemes used elsewhere exists for these fields, or, for \texttt{pred\_threat\_l2}, because the label applies only to a narrow subset of threat branches. 2 weak labels (marked with an asterisk) were evaluated through CC1 and CC2 but showed weak performance in these checks: \texttt{pred\_threat\_l1} and \texttt{species}. Two further weak labels (marked with an asterisk) were not evaluated through CC1 and CC2 due to the absence of available expert-labeled data, and were instead manually spot-checked by the authors against LLM predictions: \texttt{pred\_biome} and \texttt{pred\_efg}. All unvalidated and weakly validated labels are nonetheless included in the map for completeness. A hyphen (-) indicates that neither weak nor validation reason applies.}
\label{tab:coding_validation_status}
\end{table}
\end{landscape}
\restoregeometry

\paragraph*{CC1 coding validation:}
CC1 established a human-performance baseline for data coding, testing whether the label sets are recoverable from titles and abstracts alone; the same information constraint under which the LLM operates. For each label set, the author assigned labels to 100 records drawn from the Annotated Coding Dataset using titles and abstracts only, and compared these abstract-level labels against expert full-text labels from the Reference Coding Dataset; records where predictions defaulted to masked values (such as \texttt{Unclear} or \texttt{Not Applicable}) or fell outside the label schema were excluded from evaluation, so the effective $n$ per label set varies (see Supplementary Table~\ref{tab:cc1_coding}). All six label sets involve multi-label classification, where a record may carry multiple values per label set; this structurally depresses agreement metrics relative to single-label settings. 

CC1 achieved a macro-average weighted F1 of 84.24\% and Cohen's~$\kappa$ of 0.75 across label sets (see Supplementary Table~\ref{tab:cc1_coding}). Performance varies across label sets, reflecting differences in task complexity: the IPBES driver task (\texttt{L1}) involves only five possible categories, while IUCN threats (\texttt{L2}) spans 12 categories at level~0 alone and grows considerably as the hierarchy deepens. This difference is evident in the results: \texttt{threats\_l0} achieved an F1 of 75.06\%, while \texttt{threats\_l1}, which requires identifying the specific mechanism of a threat, fell to 57.62\%. For instance, a record covering threats from residential and commercial development may be correctly identified at \texttt{threats\_l0}, but it is often not possible to determine whether the origin is urban, commercial, or recreational expansion from an abstract alone. For \texttt{L6} (Taxonomy), lower-rank labels such as genus (F1 95.31\%) and species (F1 76.65\%) are the meaningful performance indicators; kingdom and phylum scores (99.53\% and 100\% respectively) reflect data composition: the validation sample for phylum contained a single dominant class (\textit{Chordata}) and should be interpreted conservatively. Taken together, CC1 establishes that validated labels are recoverable from titles and abstracts at the performance levels shown, providing the human baseline against which LLM performance is assessed in CC2.

\begin{table}[H]
\renewcommand{\arraystretch}{1.3}
\setlength{\tabcolsep}{4pt}
\centering
\scriptsize
\begin{tabular}{|>{\raggedright\arraybackslash}p{0.20\linewidth}
                |>{\raggedright\arraybackslash}p{0.16\linewidth}
                |>{\centering\arraybackslash}p{0.05\linewidth}
                |>{\centering\arraybackslash}p{0.11\linewidth}
                |>{\centering\arraybackslash}p{0.11\linewidth}
                |>{\centering\arraybackslash}p{0.09\linewidth}
                |>{\centering\arraybackslash}p{0.09\linewidth}|}
\hline
\textbf{Label set} & \textbf{Label} & \textbf{$n$} & \textbf{Precision (\%)} & \textbf{Recall (\%)} & \textbf{F1 (\%)} & \textbf{$\kappa$} \\
\hline
L1 -- IPBES Direct Drivers & driver & 100 & 86.60 & 79.44 & 81.94 & 0.76 \\
\hline
\multirow{2}{*}{L2 -- IUCN Threats} & threats\_l0 & 86 & 82.22 & 70.15 & 75.06 & 0.67 \\
 & threats\_l1 & 49 & 72.51 & 51.47 & 57.62 & 0.53 \\
\hline
\multirow{3}{*}{L3 -- Geographic Scope} & region & 100 & 95.02 & 90.20 & 92.38 & 0.89 \\
 & sub-region & 91 & 96.35 & 91.40 & 93.33 & 0.90 \\
 & country & 90 & 95.10 & 81.91 & 87.39 & 0.84 \\
\hline
L4 -- Ecosystem Typology & realm & 100 & 100 & 93.07 & 96.14 & 0.93 \\
\hline
L5 -- Study Attributes & study\_design & 100 & 82.31 & 76.47 & 76.76 & 0.52 \\
\hline
\multirow{6}{*}{L6 -- Taxonomy} & kingdom & 100 & 99.12 & 100 & 99.53 & 0.99 \\
 & phylum & 90 & 100 & 100 & 100 & - \\
 & class & 90 & 94.15 & 95.56 & 94.81 & 0.63 \\
 & order & 86 & 92.65 & 93.55 & 93.06 & 0.80 \\
 & genus & 78 & 96.72 & 94.25 & 95.31 & 0.94 \\
 & species & 85 & 78.48 & 76.03 & 76.65 & 0.76 \\
\hline
\end{tabular}
\caption{\textbf{CC1 coding validation results:} Across label sets, CC1 achieved an overall macro-average weighted precision of 89.21\%, weighted recall of 81.81\%, weighted F1 of 84.24\%, and weighted Cohen's~$\kappa$ of 0.75. For this summary, weighted metrics were first macro-averaged within each label set and then macro-averaged across label sets. All reported precision, recall, and F1 values are weighted metrics; $n$ denotes records valid for scoring after exclusion of masked or out-of-schema labels, such as unclear cases or instances where no ground-truth label was available. The \textit{phylum} results should be interpreted with caution because the validation data contained only one observed class (\textit{Chordata}); accordingly, Cohen's~$\kappa$ is not applicable.}
\label{tab:cc1_coding}
\end{table}


\paragraph*{CC2 coding validation:}
CC2 evaluated LLM coding labels against expert full-text labels across six label sets, providing the ground-truth accuracy check complementary to CC1 human baseline. Prompts were iteratively developed on the training split and reasoning effort configurations compared on the development split; the best-performing configuration per label set was then evaluated on the held-out test split.

On the held-out test split, CC2 achieved a macro-average weighted F1 of 83.54\% and Cohen's~$\kappa$ of 0.72 across label sets. As with CC1, all label sets involve multi-label classification, which structurally depresses agreement metrics relative to single-label settings (see Supplementary Table~\ref{tab:cc2_coding}).

Performance patterns mirror those observed in CC1 and reflect the same underlying constraints. For \texttt{L2} (IUCN Threats), \texttt{threats\_l1} (F1 57.52\%) remains considerably lower than \texttt{threats\_l0} (F1 68.43\%); as in CC1, this reflects the specificity of sub-level threat classification rather than a model limitation: the abstract-level text frequently does not contain sufficient detail to distinguish. For \texttt{L6} (Taxonomy), species-level performance (F1 57.95\%) falls notably below the CC1 human baseline (F1 76.65\%). In CC1, human annotators could leverage contextual cues, sub-species disambiguation, and manual taxonomy lookups; the LLM resolved species identity from text alone, which introduces conflation between adjacent taxa. Importantly, this does not propagate upward through the hierarchy: even where a species prediction is adjacent rather than exact: for instance, a different species within the same genus are resolved via GBIF Backbone Taxonomy lookup and remain unaffected, as confirmed by the stronger genus performance (F1 93.88\%) relative to species. Overall, CC2 broadly approaches CC1 performance across validated label sets, indicating that LLM coding from titles and abstracts is a reliable approach to large-scale systematic mapping.

\newgeometry{top=0.8cm, bottom=1.5cm, left=1cm, right=1cm}
\thispagestyle{plain}
\begin{landscape}
\begin{table}[p]
\renewcommand{\arraystretch}{1.2}
\setlength{\tabcolsep}{3pt}
\centering
\scriptsize
\begin{tabularx}{\linewidth}{|>{\raggedright\arraybackslash}p{0.07\linewidth}
                |>{\raggedright\arraybackslash}p{0.09\linewidth}
                |>{\centering\arraybackslash}X
                |>{\centering\arraybackslash}X
                |>{\centering\arraybackslash}X
                |>{\centering\arraybackslash}X
                |>{\centering\arraybackslash}X
                |>{\centering\arraybackslash}X
                |>{\centering\arraybackslash}X
                |>{\centering\arraybackslash}X
                |>{\centering\arraybackslash}X
                |>{\centering\arraybackslash}X
                |>{\centering\arraybackslash}X
                |>{\centering\arraybackslash}X
                |>{\centering\arraybackslash}X
                |>{\centering\arraybackslash}X|}
\hline
\textbf{Split} & \textbf{Metric} &
\textbf{L1 driver} & \textbf{L2 threats l0} & \textbf{L2 threats l1} &
\textbf{L3 region} & \textbf{L3 sub-region} & \textbf{L3 country} &
\textbf{L4 realm} & \textbf{L5 study design} &
\textbf{L6 kingdom} & \textbf{L6 phylum} & \textbf{L6 class} &
\textbf{L6 order} & \textbf{L6 genus} & \textbf{L6 species} \\
\hline
& Selected config & Low & Medium & High & Low & Low & Low & High & Low & High & High & High & High & High & High \\
& $n$ & 841 & 111 & 93 & 471 & 411 & 411 & 843 & 952 & 148 & 134 & 126 & 120 & 109 & 126 \\
\hline
\multirow{4}{*}{Dev Low} & Precision & 81.75 & 80.57 & 63.68 & 92.22 & 84.32 & 94.23 & 94.58 & 86.74 & 98.47 & 100 & 88.47 & 91.85 & 96.03 & 75.68 \\
& Recall & 88.96 & 69.88 & 56.12 & 89.60 & 89.50 & 75.53 & 85.31 & 82.89 & 95.97 & 98.51 & 90.62 & 88.15 & 96.55 & 69.23 \\
& F1 & 85.06 & 73.43 & 56.42 & 90.53 & 86.31 & 82.57 & 89.49 & 84.30 & 97.19 & 99.25 & 88.82 & 89.71 & 96.02 & 70.35 \\
& Cohen's $\kappa$ & 0.79 & 0.66 & 0.51 & 0.85 & 0.80 & 0.76 & 0.80 & 0.66 & 0.80 & - & 0.42 & 0.86 & 0.95 & 0.69 \\
\hline
\multirow{4}{*}{Dev Medium} & Precision & 80.92 & 82.31 & 65.27 & 90.23 & 83.83 & 94.69 & 93.93 & 86.32 & 98.21 & 100 & 87.70 & 91.91 & 95.97 & 74.56 \\
& Recall & 87.43 & 70.56 & 56.49 & 87.53 & 89.98 & 74.12 & 86.84 & 82.06 & 94.63 & 97.01 & 88.37 & 86.76 & 95.69 & 69.74 \\
& F1 & 83.79 & 73.67 & 57.58 & 88.54 & 86.18 & 82.15 & 90.02 & 83.63 & 96.34 & 98.48 & 87.35 & 89.09 & 95.51 & 70.95 \\
& Cohen's $\kappa$ & 0.78 & 0.66 & 0.51 & 0.82 & 0.80 & 0.75 & 0.81 & 0.65 & 0.75 & - & 0.29 & 0.77 & 0.94 & 0.70 \\
\hline
\multirow{4}{*}{Dev High} & Precision & 80.09 & 79.70 & 64.94 & 91.33 & 83.59 & 94.37 & 94.57 & 86.10 & 98.66 & 100 & 87.73 & 91.11 & 96.77 & 74.32 \\
& Recall & 86.67 & 70.17 & 58.97 & 89.19 & 91.41 & 74.35 & 86.84 & 81.34 & 97.32 & 98.50 & 89.84 & 88.15 & 97.39 & 70.26 \\
& F1 & 82.97 & 72.84 & 58.60 & 89.74 & 86.65 & 82.15 & 90.44 & 83.06 & 97.98 & 99.24 & 88.06 & 89.43 & 96.85 & 71.22 \\
& Cohen's $\kappa$ & 0.76 & 0.65 & 0.53 & 0.84 & 0.81 & 0.75 & 0.82 & 0.64 & 0.86 & - & 0.33 & 0.83 & 0.96 & 0.70 \\
\hline
\multirow{5}{*}{Train} & $n$ & 2515 & 325 & 295 & 1410 & 1227 & 1223 & 2523 & 2856 & 441 & 396 & 384 & 366 & 343 & 384 \\
& Precision & 80.89 & 77.76 & 69.63 & 92.25 & 86.04 & 92.20 & 96.00 & 88.36 & 98.51 & 100 & 90.47 & 90.93 & 93.92 & 61.27 \\
& Recall & 87.58 & 65.77 & 54.32 & 90.08 & 91.28 & 77.52 & 89.13 & 84.99 & 96.16 & 97.73 & 91.49 & 88.75 & 92.43 & 53.42 \\
& F1 & 84.01 & 70.84 & 59.09 & 90.76 & 88.13 & 83.57 & 92.30 & 86.37 & 97.31 & 98.85 & 90.57 & 89.36 & 92.67 & 55.03 \\
& Cohen's $\kappa$ & 0.78 & 0.63 & 0.54 & 0.86 & 0.83 & 0.78 & 0.85 & 0.70 & 0.76 & - & 0.30 & 0.75 & 0.91 & 0.54 \\
\hline
\multirow{5}{*}{Test} & $n$ & 837 & 111 & 97 & 469 & 404 & 402 & 839 & 950 & 146 & 131 & 126 & 119 & 113 & 125 \\
& Precision & 81.02 & 72.99 & 63.58 & 91.76 & 86.99 & 91.05 & 95.85 & 86.94 & 100 & 100 & 94.46 & 95.96 & 92.27 & 59.77 \\
& Recall & 87.05 & 65.73 & 56.13 & 90.02 & 91.21 & 77.08 & 89.73 & 82.97 & 97.96 & 99.24 & 96.09 & 93.55 & 95.76 & 58.29 \\
& F1 & 83.85 & 68.43 & 57.52 & 90.38 & 88.41 & 82.97 & 92.61 & 84.60 & 98.90 & 99.62 & 94.95 & 94.54 & 93.88 & 57.95 \\
& Cohen's $\kappa$ & 0.78 & 0.58 & 0.51 & 0.85 & 0.83 & 0.77 & 0.85 & 0.67 & 0.93 & - & 0.25 & 0.63 & 0.93 & 0.57 \\
\hline
\end{tabularx}
\caption{\textbf{CC2 coding validation results:} On the held-out test set, CC2 achieved an overall macro-average weighted precision of 85.41\%, weighted recall of 82.82\%, weighted F1 of 83.54\%, and weighted Cohen's~$\kappa$ of 0.72. For this summary, weighted metrics were first averaged within each label set and then macro-averaged across label sets. All reported precision, recall, and F1 values are weighted metrics; $n$ denotes records valid for scoring after exclusion of masked or out-of-schema labels, such as unclear cases or instances where no ground-truth label was available. The \textit{phylum} results should be interpreted with caution because the validation data contained only one observed class (\textit{Chordata}); accordingly, performance may be overstated and Cohen's~$\kappa$ is not applicable, reported as ``-''.}
\label{tab:cc2_coding}
\end{table}
\end{landscape}
\restoregeometry

\paragraph*{Weak labels:}
As summarized in Supplementary Table~\ref{tab:coding_validation_status}, 11 of the 23 coding labels are designated \emph{weak}, for one of three reasons. Two, \texttt{pred\_threat\_l1} and \texttt{species}, were evaluated through the full CC1 and CC2 consistency checks reported above and are weak based on those results. Two further labels, \texttt{pred\_biome} and \texttt{pred\_efg}, lacked expert-labeled ground truth and were instead audited manually. The remaining seven lacked expert-labeled data and a standardized schema against which to validate; \texttt{pred\_threat\_l2} is a partial exception, as its most specific IUCN level is structurally absent from roughly half the truth set's hierarchy branches, making validation largely inapplicable rather than merely unattempted.

\texttt{pred\_biome} and \texttt{pred\_efg} were audited by comparing the authors' manual reading of titles and abstracts against LLM predictions for a sample of 50 records each. \texttt{pred\_biome} achieved a weighted precision of 94.11\%, recall of 86.27\%, F1 of 90.01\%, and Cohen's~$\kappa$ of 0.65; \texttt{pred\_efg} achieved a weighted precision of 83.92\%, recall of 74.50\%, F1 of 76.75\%, and Cohen's~$\kappa$ of 0.63. These results are reasonably strong despite the absence of a formal validation cycle, likely because realm, biome, and ecosystem functional group are often directly inferable from the ecological description given in a title or abstract.

\newpage
\subsubsection*{Supplementary methods: Reporting checklist}
\phantomsection\label{smethods:llm_config}
A reporting checklist~\cite{feuerriegel_reporting_2026} for this systematic map is presented in Extended Data Table~\ref{tab:guide_llm_checklist}, documenting the transparency, reproducibility, and rationale behind majority of LLM usage decision in the pipeline. To further support reproducibility, the available code includes a representative subset of LLM inputs, raw model outputs, and evaluation results in their original form.



%%%%%%%%%%%%%%%%%%%%%%%%%%%%%%%%%%%%%%%%%%%%%%%%%%%%%%%%%%%%%%%
% Extended data section 
\newpage
\newpage
\section*{Extended Data}
\setcounter{table}{0}
\setcounter{figure}{0}
\renewcommand{\thetable}{\arabic{table}}
\renewcommand{\thefigure}{\arabic{figure}}
\renewcommand{\tablename}{Extended Data Table}
\renewcommand{\figurename}{Extended Data Fig.}


\newpage
% tab:guide_llm_checklist

\newgeometry{left=1.5cm, right=1.5cm, top=2cm, bottom=2cm}
\begin{landscape}
{\singlespacing
\begin{longtable}{|p{0.55\linewidth}|p{0.40\linewidth}|}

\caption{\textbf{Reporting checklist for LLM use:} The GUIDE-LLM checklist [26] is completed here to document the transparency and reproducibility of LLM use in this systematic map.}

\label{tab:guide_llm_checklist} \\
\hline
\rowcolor{green!20}
\textbf{Item} & \textbf{Answer} \\* \hline
\endfirsthead
\caption*{\textbf{Reporting checklist for LLM use (continued).}} \\
\hline
\rowcolor{green!20}
\textbf{Item} & \textbf{Answer} \\* \hline
\endhead
\hline
\multicolumn{2}{r}{\footnotesize\itshape Continued on next page\ldots} \\
\endfoot
\hline
\endlastfoot
%% SECTION A
\multicolumn{2}{|l|}{\cellcolor{green!20}\textbf{Scope of LLM use}} \\
\hline
\textbf{Item A.1:} LLMs were used in this project for: \textcolor{gray}{{\footnotesize \textbf{Explanation:} Briefly describe how and for what purposes LLMs were used in the study. This may include one or multiple stages of the research workflow. The following examples illustrate common use cases: \begin{itemize}[leftmargin=*]\setlength{\itemsep}{0pt}\setlength{\parskip}{0pt} \item \textbf{Research design} (e.g., hypothesis generation, literature search, or creating surveys/stimuli). \item \textbf{Data processing} (e.g., transcription, translation, data extraction, or data cleaning). \item \textbf{Analysis} (e.g., data labeling, summarization, pattern detection, statistical analysis, or code generation). \item \textbf{LLM as research object} (e.g., studying LLM behavior, benchmarking LLMs, or bias assessment of LLMs). \item \textbf{Participant-facing settings} (e.g., LLM used as an intervention, studying human interactions with LLM chatbots). \item \textbf{Communication} (e.g., paper writing, editing, or reviewing). \end{itemize} Depending on the specific use case described here, different checklist items may later be relevant.}} & LLM use involved eligibility screening of bibliographic records (title and abstract) using a four-question procedure (Q1--Q4) and data coding of eligible records across six label sets (L1--L6) covering anthropogenic drivers, threat classification, geographic scope, ecosystem typology, study attributes, and taxonomy. \\
\hline
\textbf{Item A.2:} Degree of automation (human-in-the-loop vs.\ fully automated): \textcolor{gray}{{\footnotesize \textbf{Explanation:} Indicate how much human oversight was involved. Specify whether each output was reviewed, edited, or approved by a person, or whether outputs were used automatically without supervision. Specify who provided oversight (e.g., student assistant, expert, PI).}} & Primarily automated. All eligibility screening and data coding decisions on individual records were made by the model, with no per-record human review during full-corpus application. Human involvement occurred at defined stages: prompt design and iterative refinement on training and development splits prior to freezing, construction of validation datasets (CC1/CC2, covering over 8,000 human-annotated coding records and approximately 2,900 screening records), evaluation of pipeline performance on a held-out test split, and a post-hoc audit of a sample of screened and coded records conducted by the authors after full-corpus application.

\\
\hline
%% SECTION B
\multicolumn{2}{|l|}{\cellcolor{green!20}\textbf{Model/system details}} \\
\hline
\textbf{Item B.1:} Model name, including provider, model size, exact version/ID, date of access, and source link (if possible): \textcolor{gray}{{\footnotesize \textbf{Explanation:} Report the exact model names (including provider, version, and date accessed). Avoid generic labels like ``ChatGPT'' or ``GPT-4''; instead, use detailed model names. If multiple models were tested, name them and briefly explain which one was used in the final study and why.}} & GPT-5-nano (version \texttt{2025-08-07}), provided by OpenAI. Accessed via the OpenAI API. The model and its version were frozen for selection prior to full-corpus application. \\
\hline
\textbf{Item B.2:} Model access (e.g., API, web interface, local) and context mode (e.g., chat mode or separate calls): \textcolor{gray}{{\footnotesize \textbf{Explanation:} Note how you accessed the models and whether you used LLMs in chat mode (ongoing conversation) or stateless mode (separate prompts). Mention the exact API name and version, since different access modes may influence responses.}} & Access via the OpenAI Responses API. Each record was processed via stateless API calls with no persistent session memory across records. For hierarchical label sets (IUCN Threats, \texttt{L2}; Ecosystem Typology, \texttt{L4}), multiple sequential stateless calls were made per record, with the output of each stage passed as context to the next. \\
\hline
\textbf{Item B.3:} Relevant LLM configurations reported (as applicable), such as temperature, max tokens, seed, and number of runs: \textcolor{gray}{{\footnotesize \textbf{Explanation:} List any configuration settings that may affect outputs, such as temperature, sampling parameters, penalties, stop sequences, number of completions or runs, quantization level, and reasoning-related settings.}} & The OpenAI Responses API was used in batch mode. Inference was stateless across records. For hierarchical classifications, later stages ran as additional stateless calls per record, with prior-stage outputs re-injected as context. Reasoning effort varied by task and was fixed before full-corpus application. Exact task-specific reasoning-effort settings, output-token limits, and any other non-default parameters are reported in the supplementary code repository. \\
\hline
\textbf{Item B.4:} Customization: \textcolor{gray}{{\footnotesize \textbf{Explanation:} Check and describe any modifications or extended capabilities incorporated into your LLM setup beyond standard inference, including fine-tuning, RAG, automated prompt optimization, web search integration, agentic workflows, or post-training refinements.}} & \begin{itemize}[leftmargin=*]\setlength{\itemsep}{0pt}\setlength{\parskip}{0pt} \item[\checked] Base model \item[\unchecked] Fine-tuning \item[\unchecked] RAG (retrieval-augmented) \item[\unchecked] Automated prompt optimization \item[\unchecked] Tool/function calling \item[\unchecked] Web search \item[\unchecked] Agentic workflows \item[\unchecked] Other adaptations \end{itemize} Description: Customization was limited to task-specific prompts, staged hierarchical inference for multi-level label sets, and a prebuilt GBIF Backbone Taxonomy cache used for taxonomic hierarchy lookup and name normalization in post-inference processing. \\
\hline
\textbf{Item B.5:} Did the LLM session(s) include persistent memory across interactions? \textcolor{gray}{{\footnotesize \textbf{Explanation:} Indicate whether the LLM could ``remember'' previous conversations (i.e., had persistent memory). Unless such memory is disabled, there may also be spillover effects from other chat windows or prior conversations.}} & \begin{itemize}[leftmargin=*]\setlength{\itemsep}{0pt}\setlength{\parskip}{0pt} \item[\unchecked] Yes \item[\checked] No \item[\unchecked] N/A \end{itemize} \\
\hline
%% SECTION C
\multicolumn{2}{|l|}{\cellcolor{green!20}\textbf{Prompts}} \\
\hline
\textbf{Item C.1:} Exact prompt(s) reported: \textcolor{gray}{{\footnotesize \textbf{Explanation:} Whenever possible, include the exact text of prompts you used. If full prompts cannot be shared, include a redacted or representative example or link to the full prompt in a repository.}} & All prompts are available in the supplementary code repository. \\
\hline
\textbf{Item C.2:} System-wide instructions (if any): \textcolor{gray}{{\footnotesize \textbf{Explanation:} Note any system-level instructions that guide the model's general behavior (e.g., ``You are a helpful assistant.''). These are commonly not directly visible but can be accessed through the API.}} & No single system-wide instruction was applied across all tasks. Each label set used a dedicated task-specific prompt. \\
\hline
%% SECTION D
\multicolumn{2}{|l|}{\cellcolor{green!20}\textbf{Data inputs \& privacy}} \\
\hline
\textbf{Item D.1:} Handling of personal or sensitive data (if any) (e.g., consent for data processing): \textcolor{gray}{{\footnotesize \textbf{Explanation:} If any personal, sensitive, or identifiable data were processed, describe how they were handled in compliance with ethical standards and data protection laws. Clarify where data were stored or processed and how applicable legal/ethical requirements were met.}} & No personal or sensitive data were processed. All model inputs consisted of titles and abstracts of bibliographic records retrieved from the Web of Science Core Collection. \\
\hline
%% SECTION E
\multicolumn{2}{|l|}{\cellcolor{green!20}\textbf{Validation \& interpretation}} \\
\hline
\textbf{Item E.1:} Human validation of LLM outputs: \textcolor{gray}{{\footnotesize \textbf{Explanation:} If relevant, describe whether and how human reviewers examined the model's outputs, and the degree of independence they had in doing so. Specify the reviewers' roles and relevant expertise. Clarify what dimensions of performance were examined and what metrics were used.}} & \begin{itemize}[leftmargin=*]\setlength{\itemsep}{0pt}\setlength{\parskip}{0pt} \item[\checked] Yes \item[\unchecked] No \item[\unchecked] N/A \end{itemize} Description: LLM labels were evaluated against the human-annotated records used in CC1 and CC2. Screening performance was assessed primarily using recall, and coding performance using precision, recall, F1, and Cohen's $\kappa$. Following full-corpus application, the authors conducted a post-hoc manual review of a sample of screened and coded records.

\\
\hline
\textbf{Item E.2:} Describe any relevant post-processing (e.g., filtering in case of format mismatches, unit conversions, etc.): \textcolor{gray}{{\footnotesize \textbf{Explanation:} Describe any steps you took to clean or reformat LLM outputs. State how you handled inconsistent or unusable outputs and whether corrections were made with an automated script or manually.}} & Post-processing was fully automated and focused on parsing, schema harmonization, and label normalization. Outputs were requested in structured JSON format; responses were parsed into standardized records, with malformed JSON repaired using a JSON-repair utility where needed, and unrecoverable failures retained with explicit error fields. For evaluation, nested outputs were flattened into tabular fields, list-like predictions were deduplicated and normalized, and known label variants were canonicalized (including threat-label aliases, geography normalization to ISO 3166-1 alpha-3, and screening decision labels). Taxa predictions were enriched post hoc with higher-rank taxonomy fields using a prebuilt GBIF lookup cache. \\
\hline
%% SECTION F
\multicolumn{2}{|l|}{\cellcolor{green!20}\textbf{Reproducibility}} \\
\hline
\textbf{Item F.1:} Code/notebooks/scripts for LLM calls shared: \textcolor{gray}{{\footnotesize \textbf{Explanation:} Indicate whether you have shared materials such as code, prompts, logs, or transcripts. Make sure sensitive information (e.g., API keys, private data) is removed. For code, make sure to add a README file.}} & \begin{itemize}[leftmargin=*]\setlength{\itemsep}{0pt}\setlength{\parskip}{0pt} \item[\checked] Yes \item[\unchecked] No \item[\unchecked] N/A \end{itemize} Link/DOI: Code repository is provided \\
\hline
%% SECTION G
\multicolumn{2}{|l|}{\cellcolor{green!20}\textbf{Competing interests}} \\
\hline
\textbf{Item G.1:} Funding, support, or other relevant relationships (including in-kind access to compute or models, or professional affiliations): \textcolor{gray}{{\footnotesize \textbf{Explanation:} Disclose any current or past funding, support, or other relevant relationships with entities that have a financial interest in LLMs. Disclose these relationships regardless of whether or not you believe they impacted the research.}} & \begin{itemize}[leftmargin=*]\setlength{\itemsep}{0pt}\setlength{\parskip}{0pt} \item[\unchecked] Yes. Description: \item[\checked] No \end{itemize} Link/DOI: \\
\hline
%% OPTIONAL ITEMS
\multicolumn{2}{|l|}{\cellcolor{yellow!30}\textbf{Optional items}} \\
\hline
\textbf{Justification for the LLM choice along the following dimensions:} \textcolor{gray}{\footnotesize \textbf{Explanation:} Briefly explain why you chose this model for your study and why it is suited to the research needs. Consider: performance, transparency, reproducibility, ethical considerations, as well as cost or ease-of-use if relevant.} & \begin{itemize}[leftmargin=*]\setlength{\itemsep}{0pt}\setlength{\parskip}{0pt} \item[\checked] Performance -- Description: GPT-5-nano was selected for formal evaluation because it is suited to large-scale classification tasks. Its performance for screening and coding was assessed through CC1 and CC2. \item[\unchecked] Transparency -- Description: \item[\unchecked] Reproducibility -- Description: \item[\unchecked] Ethical considerations -- Description: \item[\checked] Others (e.g., cost, ease-of-use) -- Description: GPT-5-nano provided a cost-effective option at the scale of this map (${\sim}2.4$ million records), enabling large-scale automated processing without significant computational cost. \end{itemize} \\
\hline
\textbf{Discussion of the rationale for the prompt design:} \textcolor{gray}{\footnotesize \textbf{Explanation:} Explain how you designed your prompts. Indicate whether you used a structured format, followed established prompt engineering guidelines, adapted prompts from earlier studies, or relied on automated prompt optimization tools. Clarify whether the design was iterative and whether few-shot examples were included.} & Prompts followed a structured, task-specific format with explicit task definitions, decision procedures, controlled label vocabularies, and output constraints. Where applicable, label definitions and decision rules were adopted verbatim from established classification schemes used in this map (see protocol). Each prompt was paired with a structured output schema enforced at the API level. Prompts were refined iteratively on training and development splits; no automated prompt optimization was used. \\
\hline
\textbf{Comparison against other methods/LLMs:} \textcolor{gray}{\footnotesize \textbf{Explanation:} If comparisons were made between models, prompts, or methods, describe how the comparison was conducted, metrics used, and fairness of inputs. If alternatives were excluded, justify why.} & \begin{itemize}[leftmargin=*]\setlength{\itemsep}{0pt} \item[\checked] Yes. Metrics: Specified in Item E.1 \item[\unchecked] No \item[\unchecked] N/A \end{itemize} \\
\hline
\textbf{Training data leakage risks addressed:} \textcolor{gray}{\footnotesize \textbf{Explanation:} Note whether risks of evaluation materials appearing in model training data were considered. Describe mitigation steps such as novelty checks or item restructuring.} & \begin{itemize}[leftmargin=*]\setlength{\itemsep}{0pt} \item[\checked] Yes: To prevent leakage within the study workflow, the held-out test split was not used for prompt development or refinement \item[\unchecked] No \item[\unchecked] N/A \end{itemize} \\
\hline

\textbf{Potential risk of bias or systematic differences in LLM behavior:} \textcolor{gray}{\footnotesize \textbf{Explanation:} Reflect on whether the model may perform differently across groups, languages, or contexts in ways that could affect conclusions. Describe any bias checks or mitigation.} & \begin{itemize}[leftmargin=*]\setlength{\itemsep}{0pt} \item[\checked] Yes. Description: Potential systematic differences may occur across geographic contexts, taxonomic naming conventions, and article phrasing styles. No formal subgroup bias audit was conducted, but consistency was partially supported through controlled vocabularies, conservative decision rules, fallback labels (\texttt{Unclear}, \texttt{Not Applicable}), schema-constrained outputs, and post hoc normalization of geography and taxonomy outputs. \item[\unchecked] No \end{itemize} \\
\hline
\textbf{Conversation transcripts:} \textcolor{gray}{\footnotesize \textbf{Explanation:} For studies involving direct researcher/participant interaction with an LLM, provide anonymized transcripts or representative examples.} & \begin{itemize}[leftmargin=*]\setlength{\itemsep}{0pt} \item[\unchecked] Yes, shared without sensitive information. Location: \item[\unchecked] No. Reason for not sharing: \item[\checked] N/A. Description: No interaction involved. \end{itemize} \\
\hline

\textbf{Relevant ethical implications of the research:} \textcolor{gray}{\footnotesize \textbf{Explanation:} Discuss broader ethical considerations beyond procedural compliance, including participant well-being, fairness, safety, autonomy, and dual-use risks.} & The main ethical implications of this work concern the use of LLMs for automated evidence synthesis. Misclassification or systematic differences in model behavior could bias downstream conclusions, particularly for brief or ambiguous abstracts. The map therefore used structured outputs, conservative rules, uncertainty labels, and evaluation against labeled data, and its outputs were interpreted in light of the reported validation results and methodological limitations. \\
\hline
\textbf{Computational resources:} \textcolor{gray}{\footnotesize \textbf{Explanation:} Report the computational resources required for the study. This may include API call counts, total tokens processed, financial costs, runtime, number and type of GPUs/CPUs used, and any other relevant compute metrics.} & API Call Count: 4.15 Million Tokens Processed: 20.27 Billion (Input = 12.96 Billion, Output = 7.32 Billion) Financial Cost: € 1,400 \\
\hline
\end{longtable}
}
\end{landscape}
\restoregeometry

\end{document}

